\documentclass[a4paper,11pt]{article}

\usepackage{fontspec}
\usepackage{newunicodechar}
\newunicodechar{π}{\ensuremath{\pi}}
\usepackage[a4paper,margin=1in]{geometry}
\usepackage{hyperref}
\hypersetup{colorlinks=true,linkcolor=blue,citecolor=blue,urlcolor=blue}
\usepackage{enumitem}
\usepackage{booktabs}
\usepackage{tabularx}
\usepackage{array}
\newcolumntype{Y}{>{\raggedright\arraybackslash}X}
\usepackage[numbers]{natbib}
\usepackage{url}

\title{Embodied AI survey}
\author{}
\date{\today}

\setlist[itemize]{noitemsep,topsep=0.25em,leftmargin=*}
\setcounter{tocdepth}{2}

\begin{document}
\maketitle

\tableofcontents
\newpage

\begin{abstract}


Recent research on Embodied AI now spans multiple methodological families, task settings, and evaluation regimes, yet those threads are still often compared under mismatched protocols \citep{Fan2025Interleave,Shao2025Large,Gundawar2025Bench,Kim2026Molmospaces,Din2025Vision,Zheng2024Survey}. Across work on problem settings and embodiment assumptions, model and policy design, data, supervision, and post-training strategy, and evaluation, safety, and deployment constraints, the same question keeps returning: which reported gains survive changes in setup, data, and evaluation practice \citep{Zhu2025Objectvla,Doshi2024Scaling,Bu2025Univla,Bai2025Towards,Zhou2025Autoeval,Dai2026Conla}. Drawing on a candidate pool of 1807 records and a 300-paper core set, the paper tracks how current approaches trade off effectiveness, generality, reproducibility, and deployment realism across representative settings \citep{Jiang2025Behavior,Eze2024Learning,Kachaev2025Mind,Xu2025Aerial,Wu2024Discrete,Lin2025Echovla}. A central finding is that many reported gains depend as much on evaluation design and experimental assumptions as on the nominal methodology itself \citep{Shi2025Diversity,Mandi2022Cacti,Wu2025Moto,Awais2023Foundational,Bai2025Embodied,Black2024Vision}. The conclusion identifies where current evidence is strongest, where protocol mismatch still obscures comparison, and which open questions should structure future work \citep{Brohan2023Vision,Bucker2024Grappa,Chen2021Learning,Chen2025Internvla,Chen2025Once,Cui2025Openhelix}.

\end{abstract}

\section{Introduction}


Recent work on Embodied AI asks systems to operate across heterogeneous tasks, embodiments, and deployment conditions that do not line up as neatly as the shared labels imply. Closely related results frequently shift the task definition, sensor access, or deployment constraint that gives a reported gain its meaning. Any synthesis that ignores those differences risks turning protocol mismatch into a false methodological ranking \citep{Fan2025Interleave,Shao2025Large,Gundawar2025Bench,Kim2026Molmospaces,Din2025Vision,Zheng2024Survey}.

The strongest disagreements in Embodied AI rarely come from a single mechanism alone. They instead track changes in metric choice, training regime, and environment assumptions that quietly alter what a result can legitimately be taken to show. The introduction therefore asks which differences actually change the meaning of a reported gain, instead of rehearsing another catalog of architectures or benchmarks \citep{Zhu2025Objectvla,Doshi2024Scaling,Bu2025Univla,Bai2025Towards,Zhou2025Autoeval,Dai2026Conla}.

The field keeps revisiting a linked set of questions about task/embodiment, policy/action design, evaluation, even when papers adopt different labels. They are more stable than method names alone because they keep the comparison attached to the mechanism, protocol, and limitation evidence that travel with each claim. Taken together, those questions keep task/embodiment, policy/action design, evaluation interpretable within one argument rather than scattering them into disconnected local summaries \citep{Jiang2025Behavior,Eze2024Learning,Kachaev2025Mind,Xu2025Aerial,Wu2024Discrete,Lin2025Echovla}.

The boundary here is comparative rather than encyclopedic. The focus stays on work around problem settings and embodiment assumptions, model and policy design, data, supervision, and post-training strategy, and evaluation, safety, and deployment constraints, because those are the places where design choices most clearly change how evidence should be read across papers. That boundary keeps the paper close to interpretation and away from a surface-level roll call of everything adjacent to the field \citep{Shi2025Diversity,Mandi2022Cacti,Wu2025Moto,Awais2023Foundational,Bai2025Embodied,Black2024Vision}.

The same interpretive problem returns in work on problem settings and embodiment assumptions, model and policy design, data, supervision, and post-training strategy, and evaluation, safety, and deployment constraints. Read across the paper, those literatures keep exposing the same dependence between technical choices and the conditions under which their evidence remains meaningful. That continuity lets later sections stay with substantive contrasts instead of re-establishing the same frame each time \citep{Brohan2023Vision,Bucker2024Grappa,Chen2021Learning,Chen2025Internvla,Chen2025Once,Cui2025Openhelix}.

The survey draws on a candidate pool of 1807 records and narrows that pool to a 300-paper core set under abstract-level evidence. The synthesis uses since 2018 as its retrieval window, so later comparisons inherit a visible scope boundary instead of an implied one. Making the evidence boundary explicit keeps later chapters focused on technical differences the cited record can actually support \citep{Dai2025Manitaskgen,Dong2026Capturing,Duggan2026Price,Fang2025From,Gu2025Safe,Hou2025Dita}.

Protocol detail is treated as part of the result rather than as background decoration. Where evidence is thinner or more heterogeneous, the synthesis uses those papers to mark boundary conditions, missing controls, or unresolved trade-offs instead of forcing a sharp ranking. The practical consequence is a stricter distinction between stable evidence and suggestive but not yet decisive comparisons \citep{Hundt2022Robots,Jiang2025Asyncvla,Jlg2025Robot,Kabir2025Socially,Kawaharazuka2025Vision,Kerr2025Robot}.

With that frame in place, the rest of the paper can follow linked technical pressures instead of slipping back into chapter-sized inventories. The useful question from this point onward is not which label sounds strongest in isolation, but which design choices keep their meaning once the same evidence discipline is applied across chapters. That perspective lets the body ask which claims remain stable under the same evidence discipline instead of merely renaming method families \citep{Khan2025Foundation,Kim2024Openvla,Kim2026Robocurate,Kunze2018Artificial,Lee2026Learning,Li2024What}.

\section{Related Work}


The field around Embodied AI is already well served by reviews that clarify local vocabularies, benchmark choices, and subarea boundaries. What is still unstable is less the existence of prior overviews than the fact that many of them combine orientation, benchmarking, and application framing in ways that make cross-paper comparison drift. The key question is therefore not whether prior surveys exist, but which comparisons they actually stabilize once their assumptions are unpacked \citep{Fan2025Interleave,Shao2025Large,Gundawar2025Bench,Kim2026Molmospaces,Din2025Vision,Zheng2024Survey}.

Surveys of problem settings and embodiment assumptions help most when they sharpen the field's language for local mechanisms and comparison targets. Those papers clarify representative tasks, systems, and local terminology, yet they often stop short of keeping task/embodiment on a common comparison scale with the rest of the field. The result is helpful local coverage without a stable basis for synthesis across neighboring parts of the field \citep{Zhu2025Objectvla,Doshi2024Scaling,Bu2025Univla,Bai2025Towards,Zhou2025Autoeval,Dai2026Conla}.

Surveys of model and policy design help most when they sharpen the field's language for local mechanisms and comparison targets. That literature is useful for orientation, yet it usually inherits the assumptions of its local subfield rather than re-reading policy/action design against neighboring evidence. Those reviews provide indispensable local structure, but they do not by themselves settle how the evidence should transfer beyond that chapter-sized lens \citep{Jiang2025Behavior,Eze2024Learning,Kachaev2025Mind,Xu2025Aerial,Wu2024Discrete,Lin2025Echovla}.

Surveys of data, supervision, and post-training strategy help most when they sharpen the field's language for local mechanisms and comparison targets. That literature is useful for orientation, yet it usually inherits the assumptions of its local subfield rather than re-reading evaluation against neighboring evidence. The result is helpful local coverage without a stable basis for synthesis across neighboring parts of the field \citep{Shi2025Diversity,Mandi2022Cacti,Wu2025Moto,Awais2023Foundational,Bai2025Embodied,Black2024Vision}.

Work on evaluation, safety, and deployment constraints is especially useful when it clarifies local task definitions, interfaces, and evaluation habits. Those papers clarify representative tasks, systems, and local terminology, yet they often stop short of keeping evaluation/deployment on a common comparison scale with the rest of the field. Those reviews provide indispensable local structure, but they do not by themselves settle how the evidence should transfer beyond that chapter-sized lens \citep{Brohan2023Vision,Bucker2024Grappa,Chen2021Learning,Chen2025Internvla,Chen2025Once,Cui2025Openhelix}.

Benchmark- and evaluation-centered papers contribute a different kind of related work. They expose metric choice, task coverage, and reproducibility as substantive methodological decisions, which is often exactly where apparent disagreements between methods become interpretable. In practice, those papers are useful less as another category label and more as calibration for what a fair comparison can mean \citep{Dai2025Manitaskgen,Dong2026Capturing,Duggan2026Price,Fang2025From,Gu2025Safe,Hou2025Dita}.

Other reviews are most helpful for scope boundaries, terminology, and historical positioning. That role is valuable even when those papers are not themselves designed for head-to-head empirical comparison. The contribution here is therefore to keep that scope-setting value while refusing to inherit their organization uncritically \citep{Hundt2022Robots,Jiang2025Asyncvla,Jlg2025Robot,Kabir2025Socially,Kawaharazuka2025Vision,Kerr2025Robot}.

What remains under-specified across these strands is how to compare evidence rather than how to enumerate topics. Existing reviews often settle one of those dimensions well while letting the others recede into background assumptions. Without that joint frame, the literature drifts back toward inventories even when the cited evidence is rich enough to support stronger synthesis \citep{Khan2025Foundation,Kim2024Openvla,Kim2026Robocurate,Kunze2018Artificial,Lee2026Learning,Li2024What}.

The value of prior surveys here is less as inherited structure than as evidence about what the field already treats as settled or unstable. That stance keeps related work in a positioning role rather than turning it into a surrogate table of contents. It also keeps the body from inheriting whichever prior review happened to be easiest to imitate \citep{Li2025Comprehensive,Li2025Embodiment,Li2026Momastage,Liu2026Affordance,Luo2026Coral,Nasiriany2026Robocasa365}.

Prior surveys mainly contribute scope boundaries, naming discipline, and a record of which trade-offs the field already foregrounds. What they do not provide by default is a shared evidential frame for judging results that move across task/embodiment, policy/action design, evaluation. That distinction is what allows related work to stay useful without taking over the argumentative structure of the paper \citep{Pan2026Scalable,Sapkota2025Vision,Shaji2026From,Sridhar2025Ricl,Tang2025Geomanip,Team2025Gemini}.

That reading of related work clears space for the main chapters to compare evidence across those parts of the field without collapsing them into one nominal method family. From that point on, the paper can shift from positioning to evidence-bearing comparison without repeating another taxonomy hand-off. Related work succeeds here only if it sharpens the later comparative discipline instead of duplicating the body chapters \citep{Torne2026Multi,Vo2025Clutter,Yadav2025Robust,Yang2025Robot,Ye2026St4Vla,Yenamandra2023Homerobot}.

\section{Problem Settings \& Embodiment}


Problem Settings \& Embodiment is where differences in task and embodiment and policy and action design stop looking like implementation detail and start changing what the evidence can support \citep{Fan2025Interleave,Shao2025Large,Gundawar2025Bench,Kim2026Molmospaces}.

Across Manipulation, navigation, and task scope, Observation, action, and embodiment interfaces, and Task diversity, environment realism, and transfer constraints, the same problem keeps returning: task and embodiment, policy and action design, and task diversity and transfer realism alter the interpretation of otherwise similar results \citep{Din2025Vision,Zheng2024Survey,Zhu2025Objectvla,Doshi2024Scaling}.

\subsection{Manipulation, navigation, and task scope}


Results in manipulation, navigation, and task scope only stay comparable when they are read against VLM, MoMA, MoMaStage, VLN, UAVs, RGB, RGB-only, RGB-D, AerialVLN, OpenFly. Differences in task setting and embodiment assumptions and observation and action interface frequently imply different evaluation setups, so the key is to compare under consistent protocols where possible. \citep{Gundawar2025Bench,Kim2026Molmospaces,Din2025Vision,Zheng2024Survey}.

Seen through task setting and embodiment assumptions, the literature separates planning-centric systems from video-conditioned systems less by labels alone than by the evidence each side stresses: extensive experiments on the AerialVLN and OpenFly benchmark validate the effectiveness of the method versus acquiring datasets that comprehensively cover diverse tasks and environments is extremely costly and difficult to scale \citep{Xu2025Aerial,Kachaev2025Mind,Dai2026Conla}. One concrete result is that comprehensive evaluation in simulation and the real world shows that Interleave-VLA offers two major benefits: improves out-of-domain generalization to unseen objects by 2x compared to text input baselines \citep{Fan2025Interleave}. These comparisons are usually reported against VLM, MoMA, MoMaStage, VLN, UAVs, RGB, RGB-only, RGB-D, AerialVLN, OpenFly \citep{Li2026Momastage,Bai2025Embodied,Kachaev2025Mind,Xu2025Aerial}.

The comparison becomes weaker once indoor mobile manipulation (MoMA) enables robots to translate natural language instructions into physical actions, yet long-horizon execution remains challenging due to cascading errors and limited generalization across diverse \citep{Li2026Momastage}. Along task setting and embodiment assumptions, papers centered on planning-centric systems emphasize that extensive experiments on the AerialVLN and OpenFly benchmark validate the effectiveness of the method, while work on benchmark-driven studies more often argues that evaluations are critical to assess progress and build better policies, but evaluation in the real world, especially at a scale that would provide statistically reliable results, is costly in terms of human time and hard to obtain \citep{Xu2025Aerial,Kachaev2025Mind,Zhou2025Autoeval}. That picture has to be read cautiously because learning-based approaches often fail to maintain logical consistency over extended horizons, while methods relying on explicit scene representations impose rigid structural assumptions that reduce adaptability in dynamic settings \citep{Li2026Momastage}.

Along task setting and embodiment assumptions, papers centered on video-conditioned systems emphasize that acquiring datasets that comprehensively cover diverse tasks and environments is extremely costly and difficult to scale, while work on benchmark-driven studies more often argues that evaluations are critical to assess progress and build better policies, but evaluation in the real world, especially at a scale that would provide statistically reliable results, is costly in terms of human time and hard to obtain \citep{Dai2026Conla,Xu2025Aerial,Zhou2025Autoeval,Kachaev2025Mind}. Representative evidence comes from reports that a benchmark suite of 8 tasks in which robots interact with the diverse scenes and richly annotated objects \citep{Kim2026Molmospaces}. The comparison becomes weaker once recent advances in vision, language, and multimodal learning have substantially accelerated progress in robotic foundation models, with robot manipulation remaining a central and challenging problem \citep{Bai2025Embodied}.

Seen through observation and action interface, the literature separates planning-centric systems from video-conditioned systems less by labels alone than by the evidence each side stresses: extensive experiments on the AerialVLN and OpenFly benchmark validate the effectiveness of the method versus acquiring datasets that comprehensively cover diverse tasks and environments is extremely costly and difficult to scale \citep{Xu2025Aerial,Kachaev2025Mind,Dai2026Conla}. For example, 26 foundational datasets, and 12 simulation platforms that collectively shape the development and evaluation of VLAs models \citep{Din2025Vision}. Seen through observation and action interface, the literature separates planning-centric systems from benchmark-driven studies less by labels alone than by the evidence each side stresses: extensive experiments on the AerialVLN and OpenFly benchmark validate the effectiveness of the method versus evaluations are critical to assess progress and build better policies, but evaluation in the real world, especially at a scale that would provide statistically reliable results, is costly in terms of human time and hard to obtain \citep{Xu2025Aerial,Kachaev2025Mind,Zhou2025Autoeval}.

The comparison becomes weaker once is a significant gap exists between benchmarks designed for high-level language instruction following, which often assume perfect low-level execution, and those for low-level robot control, which rely on simple, one-step commands \citep{Kachaev2025Mind}. A recurrent split along observation and action interface runs between video-conditioned systems and benchmark-driven studies: the former is usually supported by results showing that acquiring datasets that comprehensively cover diverse tasks and environments is extremely costly and difficult to scale, whereas the latter is more often justified by evidence that evaluations are critical to assess progress and build better policies, but evaluation in the real world, especially at a scale that would provide statistically reliable results, is costly in terms of human time and hard to obtain \citep{Dai2026Conla,Xu2025Aerial,Zhou2025Autoeval,Kachaev2025Mind}. For example, evaluations are critical to assess progress and build better policies, but evaluation in the real world, especially at a scale that would provide statistically reliable results, is costly in terms of human time and hard to obtain \citep{Zhou2025Autoeval}.

That picture has to be read cautiously because this disconnect prevents a comprehensive evaluation of integrated systems where both task planning and physical execution are critical \citep{Kachaev2025Mind}. Along generalization across tasks and environments, papers centered on planning-centric systems emphasize that extensive experiments on the AerialVLN and OpenFly benchmark validate the effectiveness of the method, while work on video-conditioned systems more often argues that acquiring datasets that comprehensively cover diverse tasks and environments is extremely costly and difficult to scale \citep{Xu2025Aerial,Kachaev2025Mind,Dai2026Conla}.

One concrete result is that extensive experiments on the AerialVLN and OpenFly benchmark validate the effectiveness of the method \citep{Xu2025Aerial}. The comparison becomes weaker once kitchen-R bridges a key gap in embodied AI research, enabling more holistic and realistic benchmarking of language-guided robotic agents \citep{Kachaev2025Mind}.

Representative evidence comes from reports that in a real-world multi-task training setting with five tasks, Discrete Policy achieves an average success rate that is 26\textbackslash{}\% higher than Diffusion Policy and 15\textbackslash{}\% higher than OpenVLA \citep{Wu2024Discrete}. The comparison becomes weaker once these requirements increase system cost and integration complexity, thus hindering practical deployment for lightweight UAVs \citep{Xu2025Aerial}.

One concrete result is that comprehensive simulated and real-world results demonstrate that EchoVLA substantially improves overall performance, e.g., it achieves the highest success rates of 0.52 on manipulation and navigation tasks and 0.31 on mobile manipulation \citep{Lin2025Echovla}. Representative evidence comes from reports that task diversity proves more critical than per-task demonstration quantity, benefiting transfer from diverse pre-training tasks to novel downstream scenarios; multi-embodiment pre-training data is o \citep{Shi2025Diversity}.

\subsection{Observation, action, and embodiment interfaces}


Results in observation, action, and embodiment interfaces only stay comparable when they are read against VLA, VLA-0, SOTA, VLM, ACoT, ACoT-VLA, EAR, IAR, LIBERO, LIBERO-Plus. Differences in policy backbone and action representation and language and vision grounding into robot actions frequently imply different evaluation setups, so the key is to compare under consistent protocols where possible. \citep{Duggan2026Price,Eze2024Learning,Fan2025Interleave,Gu2025Safe}.

Seen through policy backbone and action representation, the literature separates transformer-based policies from recent systems less by labels alone than by the evidence each side stresses: evaluations across extensive benchmarks demonstrate state-of-the-art or comparative performance in simulation versus evaluations across both simulation and real-world environments validate the superiority of MOTIF, which significantly outperforms strong baselines in few-shot transfer scenarios by 6.5\% in simulation and 43.7\% in real-world settings \citep{Zheng2025Soft,Hou2025Dita,Zhi2026Motif,Ye2026St4Vla}. The benchmark context most often includes VLA, VLA-0, SOTA, VLM, ACoT, ACoT-VLA, EAR, IAR, LIBERO, LIBERO-Plus \citep{Hou2025Dita,Zheng2025Soft,Zhong2026Acot,Torne2026Multi}. A persistent caveat is that while recent vision-language-action models trained on diverse robot datasets exhibit promising generalization capabilities with limited in-domain data, their reliance on compact action heads to predict discretized or continuous actions \citep{Hou2025Dita}.

Seen through policy backbone and action representation, the literature separates transformer-based policies from other mapped systems less by labels alone than by the evidence each side stresses: evaluations across extensive benchmarks demonstrate state-of-the-art or comparative performance in simulation versus internVLA-M1 employs a two-stage pipeline: (i) spatial grounding pre-training on over 2.3M spatial reasoning data to determine ``where to act'' by aligning instructions with visual, embodiment-agnostic positions \citep{Zheng2025Soft,Hou2025Dita,Chen2025Internvla,Sridhar2025Ricl}. For example, evaluations across both simulation and real-world environments validate the superiority of MOTIF, which significantly outperforms strong baselines in few-shot transfer scenarios by 6.5\% in simulation and 43.7\% in real-world settings \citep{Zhi2026Motif}. That picture has to be read cautiously because these intermediate reasoning are often indirect and inherently limited in their capacity to convey the full, granular information required for precise action execution \citep{Zhong2026Acot}.

Seen through policy backbone and action representation, the literature separates recent systems from other mapped systems less by labels alone than by the evidence each side stresses: evaluations across both simulation and real-world environments validate the superiority of MOTIF, which significantly outperforms strong baselines in few-shot transfer scenarios by 6.5\% in simulation and 43.7\% in real-world settings versus internVLA-M1 employs a two-stage pipeline: (i) spatial grounding pre-training on over 2.3M spatial reasoning data to determine ``where to act'' by aligning instructions with visual, embodiment-agnostic positions \citep{Zhi2026Motif,Ye2026St4Vla,Chen2025Internvla,Sridhar2025Ricl}. One concrete result is that sT4VLA achieves substantial improvements over vanilla VLA, with performance increasing from 66.1 -> 84.6 on Google Robot and from 54.7 -> 73.2 on WidowX Robot, establishing new state-of-the-art results on SimplerEnv \citep{Ye2026St4Vla}.

That picture has to be read cautiously because existing post-training approaches for VLA models are typically offline, single-robot, or task-specific, limiting effective on-policy adaptation and scalable learning from real-world interaction \citep{Pan2026Scalable}. A recurrent split along language and vision grounding into robot actions runs between transformer-based policies and recent systems: the former is usually supported by results showing that evaluations across extensive benchmarks demonstrate state-of-the-art or comparative performance in simulation, whereas the latter is more often justified by evidence that evaluations across both simulation and real-world environments validate the superiority of MOTIF, which significantly outperforms strong baselines in few-shot transfer scenarios by 6.5\% in simulation and 43.7\% in real-world settings \citep{Zheng2025Soft,Hou2025Dita,Zhi2026Motif,Ye2026St4Vla}.

One concrete result is that internVLA-M1 employs a two-stage pipeline: (i) spatial grounding pre-training on over 2.3M spatial reasoning data to determine `\texttt{where to act'' by aligning instructions with visual, embodiment-agnostic positions \citep{Chen2025Internvla}. Along language and vision grounding into robot actions, papers centered on transformer-based policies emphasize that evaluations across extensive benchmarks demonstrate state-of-the-art or comparative performance in simulation, while work on other mapped systems more often argues that internVLA-M1 employs a two-stage pipeline: (i) spatial grounding pre-training on over 2.3M spatial reasoning data to determine }`where to act'' by aligning instructions with visual, embodiment-agnostic positions \citep{Zheng2025Soft,Hou2025Dita,Chen2025Internvla,Sridhar2025Ricl}.

The comparison becomes weaker once existing cross-embodiment policies typically rely on shared-private architectures, which suffer from limited capacity of private parameters and lack explicit adaptation mechanisms \citep{Zhi2026Motif}. A recurrent split along language and vision grounding into robot actions runs between recent systems and other mapped systems: the former is usually supported by results showing that evaluations across both simulation and real-world environments validate the superiority of MOTIF, which significantly outperforms strong baselines in few-shot transfer scenarios by 6.5\% in simulation and 43.7\% in real-world settings, whereas the latter is more often justified by evidence that internVLA-M1 employs a two-stage pipeline: (i) spatial grounding pre-training on over 2.3M spatial reasoning data to determine ``where to act'' by aligning instructions with visual, embodiment-agnostic positions \citep{Zhi2026Motif,Ye2026St4Vla,Chen2025Internvla,Sridhar2025Ricl}.

Representative evidence comes from reports that 26 foundational datasets, and 12 simulation platforms that collectively shape the development and evaluation of VLAs models \citep{Din2025Vision}. The comparison becomes weaker once on Bridge V2, MVP-LAM produces more action-centric latent actions, achieving higher mutual information with ground-truth actions and improved action prediction, including under out-of-distribution evaluation \citep{Lee2026Learning}.

A recurrent split along latency and execution stability runs between transformer-based policies and recent systems: the former is usually supported by results showing that evaluations across extensive benchmarks demonstrate state-of-the-art or comparative performance in simulation, whereas the latter is more often justified by evidence that evaluations across both simulation and real-world environments validate the superiority of MOTIF, which significantly outperforms strong baselines in few-shot transfer scenarios by 6.5\% in simulation and 43.7\% in real-world settings \citep{Zheng2025Soft,Hou2025Dita,Zhi2026Motif,Ye2026St4Vla}. That picture has to be read cautiously because the results demonstrate that the LLM-driven agent can complete structured tasks and exhibits emergent adaptation and memory-guided planning, but also reveal significant limitations, such as hallucinations about the task success and poor \citep{Shaji2026From}.

That picture has to be read cautiously because these findings highlight both the potential and challenges of employing LLMs as embodied cognitive controllers for autonomous robots \citep{Shaji2026From}. Representative evidence comes from reports that effectively fine-tune OpenVLA for new settings, with especially strong generalization results in multi-task environments involving multiple objects and strong language grounding abilities \citep{Kim2024Openvla}.

\subsection{Task diversity, environment realism, and transfer constraints}


In task diversity, environment realism, and transfer constraints, the key comparison is not breadth alone but which assumptions keep a reported gain interpretable: Differences in task setting and embodiment assumptions and observation and action interface frequently imply different evaluation setups, so the key is to compare under consistent protocols where possible. \citep{Sridhar2025Ricl,Torne2026Multi,Vo2025Clutter,Wu2025Moto}. The literature diverges most when breadth versus transfer reliability: covering diverse tasks and environments improves scope, but it also exposes sharper failures when embodiments, sensors, or deployment conditions shift..

Along task setting and embodiment assumptions, papers centered on benchmark-driven studies emphasize that evaluating policies by systematically increasing task complexity and scenario perturbation to assess robustness, and quantifying performance alignment between real-world performance and its simulation counterparts, while work on control and conditioning schemes more often argues that the extensive evaluation (6k evaluation trials) shows that the approach leads to performant robotic policies and enables RT-2 to obtain a range of emergent capabilities from Internet-scale training \citep{Yang2025Robot,Tang2025Geomanip,Brohan2023Vision,Dong2026Capturing}. Representative evidence comes from reports that evaluating policies by systematically increasing task complexity and scenario perturbation to assess robustness, and quantifying performance alignment between real-world performance and its simulation counterparts \citep{Yang2025Robot}. The benchmark context most often includes roboCasa365, RoboCasa, GeoManip, AGI, VLMs, ManiTaskGen, VLAs, RCS, VLA, OpenVLA \citep{Nasiriany2026Robocasa365,Tang2025Geomanip,Dai2025Manitaskgen,Kachaev2025Mind}.

Seen through task setting and embodiment assumptions, the literature separates benchmark-driven studies from agent-style systems less by labels alone than by the evidence each side stresses: evaluating policies by systematically increasing task complexity and scenario perturbation to assess robustness, and quantifying performance alignment between real-world performance and its simulation counterparts versus while recent work has advanced such general robot policies, their training and evaluation are often limited to tasks within specific scenes, involving restricted instructions and scenarios \citep{Yang2025Robot,Tang2025Geomanip,Dai2025Manitaskgen}. Representative evidence comes from reports that the extensive evaluation (6k evaluation trials) shows that the approach leads to performant robotic policies and enables RT-2 to obtain a range of emergent capabilities from Internet-scale training \citep{Brohan2023Vision}. A recurrent split along task setting and embodiment assumptions runs between control and conditioning schemes and agent-style systems: the former is usually supported by results showing that the extensive evaluation (6k evaluation trials) shows that the approach leads to performant robotic policies and enables RT-2 to obtain a range of emergent capabilities from Internet-scale training, whereas the latter is more often justified by evidence that while recent work has advanced such general robot policies, their training and evaluation are often limited to tasks within specific scenes, involving restricted instructions and scenarios \citep{Brohan2023Vision,Dong2026Capturing,Dai2025Manitaskgen}.

One concrete result is that evaluations across both simulation and real-world environments validate the superiority of MOTIF, which significantly outperforms strong baselines in few-shot transfer scenarios by 6.5\% in simulation and 43.7\% in real-world settings \citep{Zhi2026Motif}. That picture has to be read cautiously because geoManip learns in-context and provides five appealing human-robot interaction features: on-the-fly policy adaptation, learning from human demonstrations, learning from failure cases, long-horizon action planning \citep{Tang2025Geomanip}. A recurrent split along observation and action interface runs between benchmark-driven studies and control and conditioning schemes: the former is usually supported by results showing that evaluating policies by systematically increasing task complexity and scenario perturbation to assess robustness, and quantifying performance alignment between real-world performance and its simulation counterparts, whereas the latter is more often justified by evidence that the extensive evaluation (6k evaluation trials) shows that the approach leads to performant robotic policies and enables RT-2 to obtain a range of emergent capabilities from Internet-scale training \citep{Yang2025Robot,Tang2025Geomanip,Brohan2023Vision,Dong2026Capturing}.

Representative evidence comes from reports that dataset scale, and environment variation on generalization \citep{Nasiriany2026Robocasa365}. The comparison becomes weaker once extensive evaluations on both simulations and real-world scenarios demonstrate GeoManip's state-of-the-art performance, with superior out-of-distribution generalization while avoiding costly model training \citep{Tang2025Geomanip}.

Along observation and action interface, papers centered on benchmark-driven studies emphasize that evaluating policies by systematically increasing task complexity and scenario perturbation to assess robustness, and quantifying performance alignment between real-world performance and its simulation counterparts, while work on agent-style systems more often argues that while recent work has advanced such general robot policies, their training and evaluation are often limited to tasks within specific scenes, involving restricted instructions and scenarios \citep{Yang2025Robot,Tang2025Geomanip,Dai2025Manitaskgen}. Representative evidence comes from reports that while recent work has advanced such general robot policies, their training and evaluation are often limited to tasks within specific scenes, involving restricted instructions and scenarios \citep{Dai2025Manitaskgen}.

The comparison becomes weaker once while recent work has advanced such general robot policies, their training and evaluation are often limited to tasks within specific scenes, involving restricted instructions and scenarios \citep{Dai2025Manitaskgen}. Along observation and action interface, papers centered on control and conditioning schemes emphasize that the extensive evaluation (6k evaluation trials) shows that the approach leads to performant robotic policies and enables RT-2 to obtain a range of emergent capabilities from Internet-scale training, while work on agent-style systems more often argues that while recent work has advanced such general robot policies, their training and evaluation are often limited to tasks within specific scenes, involving restricted instructions and scenarios \citep{Brohan2023Vision,Dong2026Capturing,Dai2025Manitaskgen}.

A persistent caveat is that existing benchmarks also typically rely on manual annotation of limited tasks in a few scenes \citep{Dai2025Manitaskgen}. A recurrent split along generalization across tasks and environments runs between benchmark-driven studies and control and conditioning schemes: the former is usually supported by results showing that evaluating policies by systematically increasing task complexity and scenario perturbation to assess robustness, and quantifying performance alignment between real-world performance and its simulation counterparts, whereas the latter is more often justified by evidence that the extensive evaluation (6k evaluation trials) shows that the approach leads to performant robotic policies and enables RT-2 to obtain a range of emergent capabilities from Internet-scale training \citep{Yang2025Robot,Tang2025Geomanip,Brohan2023Vision,Dong2026Capturing}.

A persistent caveat is that is a significant gap exists between benchmarks designed for high-level language instruction following, which often assume perfect low-level execution, and those for low-level robot control, which rely on simple, one-step commands \citep{Kachaev2025Mind}. That picture has to be read cautiously because this disconnect prevents a comprehensive evaluation of integrated systems where both task planning and physical execution are critical \citep{Kachaev2025Mind}.

For example, task diversity proves more critical than per-task demonstration quantity, benefiting transfer from diverse pre-training tasks to novel downstream scenarios; multi-embodiment pre-training data is o \citep{Shi2025Diversity}. One concrete result is that in a simulated kitchen environment, CACTI trains a single policy to perform 18 semantic tasks across 100 layout variations for each individual task \citep{Mandi2022Cacti}.

\section{Model \& Policy Architectures}


In Model \& Policy Architectures, the central question is how policy and action design and planning and world-model role change the meaning of reported progress \citep{Shao2025Large,Gu2025Usim,Li2025Survey,Xu2025Anatomy}.

Vision-language-action and policy backbones, World models, planning, and reasoning, and Transfer across embodiments and viewpoints each expose that pressure from a different angle, showing why policy and action design, planning and world-model role, and evaluation cannot be treated as interchangeable background choices \citep{Zheng2025Soft,Yu2025Survey,Zhang2025Pure,Yadav2025Robust}.

\subsection{Vision-language-action and policy backbones}


One recurrent split in vision-language-action and policy backbones runs between planning-centric systems and transformer-based policies along policy backbone and action representation. A tension around policy backbone and action representation and language and vision grounding into robot actions, motivating a protocol-aware synthesis rather than per-paper summaries. \citep{Lisondra2025Embodied,Liu2026World,Long2025Survey,Mai2024From}.

Seen through policy backbone and action representation, the literature separates planning-centric systems from transformer-based policies less by labels alone than by the evidence each side stresses: extensive evaluations across natural language processing benchmarks, robotic simulation environments, and real-world experiments confirm the superiority of MoTVLA in both fast-slow reasoning and manipulation task performance versus evaluations across extensive benchmarks demonstrate state-of-the-art or comparative performance in simulation \citep{Huang2025Motvla,Zheng2025Soft,Hou2025Dita}. The benchmark context most often includes VLA, VLMs, VLM, MoTVLA, VLAPS, MCTS, VLA-derived, VLA-only, VLA-0, SOTA \citep{Huang2025Motvla,Neary2025Improving,Hou2025Dita,Zheng2025Soft}.

A persistent caveat is that despite promising advances, existing approaches face two challenges: limited language steerability when no generated reasoning is used as a condition, or significant inference latency when reasoning is incorporated \citep{Huang2025Motvla}. That picture has to be read cautiously because pre-trained vision-language-action (VLA) models offer a promising foundation for generalist robot policies, but often produce brittle behaviors or unsafe failures when deployed zero-shot in out-of-distribution scenarios \citep{Neary2025Improving}.

One concrete result is that categorizing current techniques into three core pillars: Efficient Model Design, focusing on efficient architectures and model compression; Efficient Training \citep{Yu2025Survey}. The comparison becomes weaker once while recent vision-language-action models trained on diverse robot datasets exhibit promising generalization capabilities with limited in-domain data, their reliance on compact action heads to predict discretized or continuous actions \citep{Hou2025Dita}.

Seen through language and vision grounding into robot actions, the literature separates planning-centric systems from transformer-based policies less by labels alone than by the evidence each side stresses: extensive evaluations across natural language processing benchmarks, robotic simulation environments, and real-world experiments confirm the superiority of MoTVLA in both fast-slow reasoning and manipulation task performance versus evaluations across extensive benchmarks demonstrate state-of-the-art or comparative performance in simulation \citep{Huang2025Motvla,Zheng2025Soft,Hou2025Dita}. That picture has to be read cautiously because despite their success in controlled environments, reliable real-world deployment is severely hindered by their fragility to visual disturbances \citep{Orjuela2026Improving}.

One concrete result is that a simulation-based multi-task Vision-Language-Action (VLA) dataset for underwater robots \citep{Gu2025Usim}. That picture has to be read cautiously because robotic manipulation, a key frontier in robotics and embodied AI, requires precise motor control and multimodal understanding, yet traditional rule-based methods fail to scale or generalize in unstructured, novel environments \citep{Shao2025Large}.

Seen through latency and execution stability, the literature separates planning-centric systems from transformer-based policies less by labels alone than by the evidence each side stresses: extensive evaluations across natural language processing benchmarks, robotic simulation environments, and real-world experiments confirm the superiority of MoTVLA in both fast-slow reasoning and manipulation task performance versus evaluations across extensive benchmarks demonstrate state-of-the-art or comparative performance in simulation \citep{Huang2025Motvla,Zheng2025Soft,Hou2025Dita}.

One concrete result is that the main contribution is a detailed breakdown of the five biggest challenges in: Representation, Execution, Generalization, Safety, and Dataset and Evaluation \citep{Xu2025Anatomy}.

A persistent caveat is that trace the history through key Milestones, and then dive deep into the core Challenges that define recent research frontier \citep{Xu2025Anatomy}.

One concrete result is that categorizing current techniques into three core pillars: Efficient Model Design, focusing on efficient architectures and model compression; Efficient Training, w \citep{Yu2025Survey}.

\subsection{World models, planning, and reasoning}


Results in world models, planning, and reasoning only stay comparable when they are read against VLA, VLMs, VLM, MLLMs, FPS, OOD, DreamDojo, VLA-Loop, SANS, UWM. For World models, planning, and reasoning, planning and world-model role and long-horizon state handling is a recurring axis of variation, and results are easiest to interpret when protocols and failure assumptions are explicit. \citep{Chen2026Accelerating,Cui2026Contact,Feng2025Survey,Fu2025Metis}.

The benchmark context most often includes VLA, VLMs, VLM, MLLMs, FPS, OOD, DreamDojo, VLA-Loop, SANS, UWM \citep{Xu2026Underwater,Bai2025Embodied,Tang2025Mind,Zheng2025Multimodal}. A persistent caveat is that yet: despite rapid advances in learning-based planning and control, reliable autonomy in real ocean environments remains fundamentally constrained by tightly coupled physical limits \citep{Xu2026Underwater}.

A recurrent split along planning and world-model role runs between planning-centric systems and agent-style systems: the former is usually supported by results showing that categorizing recent progress in multimodal large language models (MLLMs) and introducing open benchmarks for evaluation, whereas the latter is more often justified by evidence that while large language models have become the prevailing approach for agentic reasoning and planning, their success in symbolic domains does not readily translate to the physical world \citep{Tang2025Mind,Zheng2025Multimodal,Felicia2026From}. A persistent caveat is that hydrodynamic uncertainty, partial observability, bandwidth-limited communication, and energy scarcity are not independent challenges; they interact within the closed perception-planning-control loop and often amplify one another over time \citep{Xu2026Underwater}.

That picture has to be read cautiously because recent advances in vision, language, and multimodal learning have substantially accelerated progress in robotic foundation models, with robot manipulation remaining a central and challenging problem \citep{Bai2025Embodied}. A recurrent split along long-horizon state handling runs between planning-centric systems and agent-style systems: the former is usually supported by results showing that categorizing recent progress in multimodal large language models (MLLMs) and introducing open benchmarks for evaluation, whereas the latter is more often justified by evidence that while large language models have become the prevailing approach for agentic reasoning and planning, their success in symbolic domains does not readily translate to the physical world \citep{Tang2025Mind,Zheng2025Multimodal,Felicia2026From}.

The comparison becomes weaker once while internet-scale data has enabled broad reasoning capabilities in AI systems, grounding these abilities in physical action remains a major challenge \citep{Tang2025Mind}.

A persistent caveat is that systematic reviews and publicly available benchmarks for these models remain limited \citep{Zheng2025Multimodal}.

One concrete result is that on the CALVIN robot manipulation benchmark, DeeR demonstrates significant reductions in computational costs of LLM by 5.2-6.5x and GPU memory of LLM by 2-6x without compromising performance \citep{Yue2024Deer}.

The comparison becomes weaker once modeling these world dynamics, especially for dexterous robotics tasks, poses significant challenges due to limited data coverage and scarce action labels \citep{Gao2026Dreamdojo}.

That picture has to be read cautiously because systematic evaluation on multiple challenging out-of-distribution (OOD) benchmarks verifies the significance of the method for simulating open-world, contact-rich tasks, paving the way for general-purpose robot world models \citep{Gao2026Dreamdojo}.

For example, categorizing recent progress in multimodal large language models (MLLMs) and introducing open benchmarks for evaluation \citep{Zheng2025Multimodal}.

\subsection{Transfer across embodiments and viewpoints}


Results in transfer across embodiments and viewpoints only stay comparable when they are read against DRL, MDPs, DDPG, TD3, PPO, SAC, LLM-based, VLA, RT-2, HIL. A tension around transfer target (embodiment and viewpoint and environment) and adaptation mechanism and invariance, motivating a protocol-aware synthesis rather than per-paper summaries. \citep{Cheang2025Technical,Chen2025Internvla,Chen2025Survey,Da2025Survey}.

A recurrent split along transfer target (embodiment and viewpoint and environment) runs between control and conditioning schemes and agent-style systems: the former is usually supported by results showing that the extensive evaluation (6k evaluation trials) shows that the approach leads to performant robotic policies and enables RT-2 to obtain a range of emergent capabilities from Internet-scale training, whereas the latter is more often justified by evidence that reliance on human supervisors imposes a 1:1 supervision ratio that limits scalability, suffers from operator fatigue over extended sessions, and introduces high variance due to inconsistent human proficiency \citep{Brohan2023Vision,Chen2026Accelerating}. One concrete result is that the extensive evaluation (6k evaluation trials) shows that the approach leads to performant robotic policies and enables RT-2 to obtain a range of emergent capabilities from Internet-scale training \citep{Brohan2023Vision}.

These comparisons are usually reported against DRL, MDPs, DDPG, TD3, PPO, SAC, LLM-based, VLA, RT-2, HIL \citep{Ter2025Taxonomy,Seo2025Equivariant,Brohan2023Vision,Chen2026Accelerating}. A persistent caveat is that reinforcement learning (RL) has become a foundational approach for enabling intelligent robotic behavior in dynamic and uncertain environments \citep{Ter2025Taxonomy}.

Representative evidence comes from reports that reliance on human supervisors imposes a 1:1 supervision ratio that limits scalability, suffers from operator fatigue over extended sessions, and introduces high variance due to inconsistent human proficiency \citep{Chen2026Accelerating}. A persistent caveat is that equivariant neural networks overcome these limitations by explicitly integrating symmetry and invariance into their architectures, leading to improved efficiency and generalization \citep{Seo2025Equivariant}.

A recurrent split along adaptation mechanism and invariance runs between control and conditioning schemes and agent-style systems: the former is usually supported by results showing that the extensive evaluation (6k evaluation trials) shows that the approach leads to performant robotic policies and enables RT-2 to obtain a range of emergent capabilities from Internet-scale training, whereas the latter is more often justified by evidence that reliance on human supervisors imposes a 1:1 supervision ratio that limits scalability, suffers from operator fatigue over extended sessions, and introduces high variance due to inconsistent human proficiency \citep{Brohan2023Vision,Chen2026Accelerating}. Representative evidence comes from reports that embodiment mixture: naively pooling heterogeneous robot datasets often induces negative transfer rather than gains, underscoring the fragility of indiscriminate data scaling \citep{Wang2026Rethinking}.

The comparison becomes weaker once reliance on human supervisors imposes a 1:1 supervision ratio that limits scalability, suffers from operator fatigue over extended sessions, and introduces high variance due to inconsistent human proficiency \citep{Chen2026Accelerating}. Representative evidence comes from reports that this paradigm shift addresses intrinsic limitations in classical Reinforcement Learning (RL), particularly the limited expressivity of standard unimodal policy distributions in capturing complex, multi-modal behaviors embedded in diverse \citep{Shao2025Generative}.

The comparison becomes weaker once a fundamental tension limits this approach: language is often too abstract to guide the concrete physical understanding required for robust manipulation \citep{Cui2026Contact}. For example, evaluations across both simulation and real-world environments validate the superiority of MOTIF, which significantly outperforms strong baselines in few-shot transfer scenarios by 6.5\% in simulation and 43.7\% in real-world settings \citep{Zhi2026Motif}.

Along cross-embodiment generalization limits, papers centered on control and conditioning schemes emphasize that the extensive evaluation (6k evaluation trials) shows that the approach leads to performant robotic policies and enables RT-2 to obtain a range of emergent capabilities from Internet-scale training, while work on agent-style systems more often argues that reliance on human supervisors imposes a 1:1 supervision ratio that limits scalability, suffers from operator fatigue over extended sessions, and introduces high variance due to inconsistent human proficiency \citep{Brohan2023Vision,Chen2026Accelerating}.

That picture has to be read cautiously because embodiment mixture: naively pooling heterogeneous robot datasets often induces negative transfer rather than gains, underscoring the fragility of indiscriminate data scaling \citep{Wang2026Rethinking}.

One concrete result is that the method outperforms strong baselines, with a 16\% gain on the SIMPLER benchmark and a 13\% improvement across real world manipulation tasks \citep{Routray2025Vipra}.

\section{Data, Training \& Post-Training}


Data, Training \& Post-Training is where differences in evaluation and data scaling and shift stop looking like implementation detail and start changing what the evidence can support \citep{Mu2024Adademo,Hejna2024Optimizing,Lin2026Systematic,Wang2026Rethinking}.

Pretraining data and supervision, Post-training, feedback, and continual improvement, and Data quality, scaling, and distribution shift each expose that pressure from a different angle, showing why evaluation and data scaling and shift cannot be treated as interchangeable background choices \citep{Zhu2025Unified,Cui2026Contact,Dass2025Datamil,Lee2025Class}.

\subsection{Pretraining data and supervision}


Results in pretraining data and supervision only stay comparable when they are read against FPS, OOD, DreamDojo, VLA, VQ-VAE, ConLA, UWM, VLAs, RCS, OpenVLA. A tension around data source and supervision and pretraining versus post-training, motivating a protocol-aware synthesis rather than per-paper summaries. \citep{Mazza2026Improving,Mei2026Video,NVIDIA2025Gr00T,Nasiriany2026Robocasa365}.

The benchmark context most often includes FPS, OOD, DreamDojo, VLA, VQ-VAE, ConLA, UWM, VLAs, RCS, OpenVLA \citep{Gao2026Dreamdojo,Dai2026Conla,Li2025Mask2Iv,Zhu2025Unified}. The comparison becomes weaker once modeling these world dynamics, especially for dexterous robotics tasks, poses significant challenges due to limited data coverage and scarce action labels \citep{Gao2026Dreamdojo}.

That picture has to be read cautiously because systematic evaluation on multiple challenging out-of-distribution (OOD) benchmarks verifies the significance of the method for simulating open-world, contact-rich tasks, paving the way for general-purpose robot world models \citep{Gao2026Dreamdojo}. Seen through data source and supervision, the literature separates control and conditioning schemes from other mapped systems less by labels alone than by the evidence each side stresses: the experiments also provide an extensive evaluation of Octo, OpenVLA, and Pi Zero on multiple robots and shed light on how simulation data can improve real-world policy performance versus learning \textbackslash{}emph\{latent actions\} from diverse human videos enables scaling robot learning beyond embodiment-specific robot datasets, and these latent actions have recently been used as pseudo-action labels for vision-language-action (VLA) \citep{Jlg2025Robot,Lee2026Learning,Lin2026Systematic}.

For example, embodiment mixture: naively pooling heterogeneous robot datasets often induces negative transfer rather than gains, underscoring the fragility of indiscriminate data scaling \citep{Wang2026Rethinking}. The comparison becomes weaker once acquiring datasets that comprehensively cover diverse tasks and environments is extremely costly and difficult to scale \citep{Dai2026Conla}.

The comparison becomes weaker once in contrast, human demonstration videos offer a rich and scalable source of diverse scenes and manipulation behaviors, yet their lack of explicit action supervision hinders direct utilization \citep{Dai2026Conla}. Representative evidence comes from reports that explicitly conditioning action generation on chain-of-thought traces learned from co-training data does not improve performance in the simulation benchmark \citep{Lin2026Systematic}.

That picture has to be read cautiously because nevertheless: since the training objective primarily focuses on reconstructing visual appearances rather than capturing inter-frame dynamics, the learned representations tend to rely on spurious visual cues, leading to shortcut learning \citep{Dai2026Conla}. A recurrent split along pretraining vs. post-training runs between control and conditioning schemes and other mapped systems: the former is usually supported by results showing that the experiments also provide an extensive evaluation of Octo, OpenVLA, and Pi Zero on multiple robots and shed light on how simulation data can improve real-world policy performance, whereas the latter is more often justified by evidence that learning \textbackslash{}emph\{latent actions\} from diverse human videos enables scaling robot learning beyond embodiment-specific robot datasets, and these latent actions have recently been used as pseudo-action labels for vision-language-action (VLA) \citep{Jlg2025Robot,Lee2026Learning,Lin2026Systematic}.

The comparison becomes weaker once while recent studies show that masks can serve as effective control signals and enhance generation quality, obtaining dense and precise mask annotations remains a major challenge for real-world use \citep{Li2025Mask2Iv}.

The comparison becomes weaker once scaling imitation learning for large robot foundation models remains challenging due to its reliance on high-quality expert demonstrations \citep{Zhu2025Unified}.

A persistent caveat is that yet: despite rapid advances in learning-based planning and control, reliable autonomy in real ocean environments remains fundamentally constrained by tightly coupled physical limits \citep{Xu2026Underwater}.

Representative evidence comes from reports that the experiments also provide an extensive evaluation of Octo, OpenVLA, and Pi Zero on multiple robots and shed light on how simulation data can improve real-world policy performance \citep{Jlg2025Robot}.

\subsection{Post-training, feedback, and continual improvement}


The strongest results in post-training, feedback, and continual improvement do not speak for themselves. They have to be read through the trade-off that structures the subsection: Differences in data source and supervision and pretraining versus post-training frequently imply different evaluation setups, so the key is to compare under consistent protocols where possible. \citep{Pan2026Scalable,Ai2025Towards,Bai2025Embodied,Bai2025Towards}. That trade-off sharpens once coverage versus depth: broader surveys capture more of the landscape but risk shallow treatment of individual methods, while focused analyses miss the wider context..

A recurrent split along data source and supervision runs between benchmark-driven studies and control and conditioning schemes: the former is usually supported by results showing that while recent work has advanced such general robot policies, their training and evaluation are often limited to tasks within specific scenes, involving restricted instructions and scenarios, whereas the latter is more often justified by evidence that where it shows significant promise in long-horizon manipulation, data efficiency in post-training, and strong generalizability to novel configurations \citep{Dai2025Manitaskgen,Jain2025Polaris,Li2026Causal,Jlg2025Robot}. For example, it comprises RoboInter-Tool, a lightweight GUI that enables semi-automatic annotation of diverse representations, and RoboInter-Data, a large-scale dataset containing over 230k episodes across 571 diverse scenes, which provides dense \citep{Li2026Robointer}. These comparisons are usually reported against roboCasa365, RoboCasa, AGI, VLMs, ManiTaskGen, PolaRiS, LingBot-VA, VLAs, RCS, VLA \citep{Nasiriany2026Robocasa365,Dai2025Manitaskgen,Jain2025Polaris,Li2026Causal}.

Seen through data source and supervision, the literature separates benchmark-driven studies from planning-centric systems less by labels alone than by the evidence each side stresses: while recent work has advanced such general robot policies, their training and evaluation are often limited to tasks within specific scenes, involving restricted instructions and scenarios versus encompassing cost-effective data generation and action prediction in imitation learning, dynamics and rewards modeling in reinforcement \citep{Dai2025Manitaskgen,Jain2025Polaris,Mei2026Video}. For example, as well as the autonomous dataset of 30.5K trajectories collected across five tabletop environments \citep{Zhou2024Autonomous}. A persistent caveat is that while recent work has advanced such general robot policies, their training and evaluation are often limited to tasks within specific scenes, involving restricted instructions and scenarios \citep{Dai2025Manitaskgen}.

Along data source and supervision, papers centered on control and conditioning schemes emphasize that where it shows significant promise in long-horizon manipulation, data efficiency in post-training, and strong generalizability to novel configurations, while work on planning-centric systems more often argues that encompassing cost-effective data generation and action prediction in imitation learning, dynamics and rewards modeling in reinforcement \citep{Li2026Causal,Jlg2025Robot,Mei2026Video}. Representative evidence comes from reports that 000 simulation rollouts and 2,835 real-world rollouts \citep{Lin2026Systematic}. A persistent caveat is that existing benchmarks also typically rely on manual annotation of limited tasks in a few scenes \citep{Dai2025Manitaskgen}.

Along pretraining vs. post-training, papers centered on benchmark-driven studies emphasize that while recent work has advanced such general robot policies, their training and evaluation are often limited to tasks within specific scenes, involving restricted instructions and scenarios, while work on control and conditioning schemes more often argues that where it shows significant promise in long-horizon manipulation, data efficiency in post-training, and strong generalizability to novel configurations \citep{Dai2025Manitaskgen,Jain2025Polaris,Li2026Causal,Jlg2025Robot}. For example, evaluation in simulation offers a scalable complement to real world evaluations, but the visual and physical domain gap between existing simulation benchmarks and the real world has made them an unreliable signal for policy improvement \citep{Jain2025Polaris}. That picture has to be read cautiously because a significant challenge for robot learning research is the ability to accurately measure and compare the performance of robot policies \citep{Jain2025Polaris}.

Along pretraining vs. post-training, papers centered on benchmark-driven studies emphasize that while recent work has advanced such general robot policies, their training and evaluation are often limited to tasks within specific scenes, involving restricted instructions and scenarios, while work on planning-centric systems more often argues that encompassing cost-effective data generation and action prediction in imitation learning, dynamics and rewards modeling in reinforcement \citep{Dai2025Manitaskgen,Jain2025Polaris,Mei2026Video}. The comparison becomes weaker once benchmarking in robotics is historically challenging due to the stochasticity, reproducibility, and time-consuming nature of real-world rollouts \citep{Jain2025Polaris}.

Seen through pretraining vs. post-training, the literature separates control and conditioning schemes from planning-centric systems less by labels alone than by the evidence each side stresses: where it shows significant promise in long-horizon manipulation, data efficiency in post-training, and strong generalizability to novel configurations versus encompassing cost-effective data generation and action prediction in imitation learning, dynamics and rewards modeling in reinforcement \citep{Li2026Causal,Jlg2025Robot,Mei2026Video}. One concrete result is that the experiments also provide an extensive evaluation of Octo, OpenVLA, and Pi Zero on multiple robots and shed light on how simulation data can improve real-world policy performance \citep{Jlg2025Robot}.

The comparison becomes weaker once this challenge is exacerbated for recent generalist policies, which has to be evaluated across a wide variety of scenes and tasks \citep{Jain2025Polaris}. A recurrent split along adaptation and transfer across embodiments runs between benchmark-driven studies and control and conditioning schemes: the former is usually supported by results showing that evaluation in simulation offers a scalable complement to real world evaluations, but the visual and physical domain gap between existing simulation benchmarks and the real world has made them an unreliable signal for policy improvement, whereas the latter is more often justified by evidence that the experiments also provide an extensive evaluation of Octo, OpenVLA, and Pi Zero on multiple robots and shed light on how simulation data can improve real-world policy performance \citep{Jain2025Polaris,Dai2025Manitaskgen,Jlg2025Robot,Li2026Causal}.

The comparison becomes weaker once evaluation in simulation offers a scalable complement to real world evaluations, but the visual and physical domain gap between existing simulation benchmarks and the real world has made them an unreliable signal for policy improvement \citep{Jain2025Polaris}. Representative evidence comes from reports that while recent work has advanced such general robot policies, their training and evaluation are often limited to tasks within specific scenes, involving restricted instructions and scenarios \citep{Dai2025Manitaskgen}.

For example, is the majority of current robot learning datasets and benchmarks mainly focus on stationary robot arms, and the few existing humanoid datasets are either confined to fixed environments or limited in task diversity, often lacking \citep{Zhao2025Humanoid}. For example, general visual representations learned from web-scale datasets for robotics have achieved great success in recent years, enabling data-efficient robot learning on manipulation tasks; yet these pre-trained representations are mostly on 2D \citep{Hou2025Visual}.

\subsection{Data quality, scaling, and distribution shift}


What matters in data quality, scaling, and distribution shift is less the headline result than the conditions that make it travel: Data quality, scaling, and distribution shift methods emphasize data quality and coverage bias and scaling efficiency versus data fidelity trade-offs, but synthesis is clearest when claims are tied to explicit evaluation settings and reporting conventions. \citep{Hou2025Dita,Zhang2025Robowheel,Mazza2026Improving,Spiridonov2025Generalist}. That is where the literature starts to split, because scale versus data fidelity: larger embodied corpora improve coverage, but noisy supervision, heterogeneous sources, and distribution shift can erode downstream robustness..

Seen through data quality and coverage bias, the literature separates video-conditioned systems from recent systems less by labels alone than by the evidence each side stresses: the approach has several advantages: it enjoys the benefit of real data to minimize the sim-to-real gap; it leverages the scalability of simulation; and it can generalize a pretrained VLA to a target domain with fully automated versus embodiment mixture: naively pooling heterogeneous robot datasets often induces negative transfer rather than gains, underscoring the fragility of indiscriminate data scaling \citep{Fang2025Rebot,Chen2021Learning,Wang2026Rethinking,Luo2026Being}. For example, the approach has several advantages: it enjoys the benefit of real data to minimize the sim-to-real gap; it leverages the scalability of simulation; and it can generalize a pretrained VLA to a target domain with fully automated \citep{Fang2025Rebot}.

These comparisons are usually reported against VLA, VLAs, ReBot, SimplerEnv, OpenVLA, RGB, DVD, YouTube, WidowX200, ACT \citep{Fang2025Rebot,Chen2021Learning,Mazza2026Improving,Wang2026Rethinking}. A persistent caveat is that is the high cost of real-world data collection hinders further data scaling, thereby restricting the generalizability of VLAs \citep{Fang2025Rebot}.

A recurrent split along data quality and coverage bias runs between video-conditioned systems and earlier related studies: the former is usually supported by results showing that the approach has several advantages: it enjoys the benefit of real data to minimize the sim-to-real gap; it leverages the scalability of simulation; and it can generalize a pretrained VLA to a target domain with fully automated, whereas the latter is more often justified by evidence that while deep learning on large and diverse datasets has shown promise as a path towards such generalization in computer vision and natural language, collecting high quality datasets of robotic interaction at scale remains an open challenge \citep{Fang2025Rebot,Chen2021Learning,Yu2023Scaling}. Representative evidence comes from reports that embodiment mixture: naively pooling heterogeneous robot datasets often induces negative transfer rather than gains, underscoring the fragility of indiscriminate data scaling \citep{Wang2026Rethinking}.

The comparison becomes weaker once a novel real-to-sim-to-real approach for scaling real robot datasets and adapting VLA models to target domains, which is the last-mile deployment challenge in robot manipulation \citep{Fang2025Rebot}. Seen through data quality and coverage bias, the literature separates recent systems from earlier related studies less by labels alone than by the evidence each side stresses: embodiment mixture: naively pooling heterogeneous robot datasets often induces negative transfer rather than gains, underscoring the fragility of indiscriminate data scaling versus while deep learning on large and diverse datasets has shown promise as a path towards such generalization in computer vision and natural language, collecting high quality datasets of robotic interaction at scale remains an open challenge \citep{Wang2026Rethinking,Luo2026Being,Chen2021Learning,Yu2023Scaling}.

That picture has to be read cautiously because reBot replays real-world robot trajectories in simulation to diversify manipulated objects (real-to-sim), and integrates the simulated movements with inpainted real-world background to synthesize physically realistic and temporally \citep{Fang2025Rebot}. That picture has to be read cautiously because the approach has several advantages: it enjoys the benefit of real data to minimize the sim-to-real gap; it leverages the scalability of simulation; and it can generalize a pretrained VLA to a target domain with fully automated \citep{Fang2025Rebot}.

Seen through scaling efficiency vs. data fidelity, the literature separates video-conditioned systems from earlier related studies less by labels alone than by the evidence each side stresses: the approach has several advantages: it enjoys the benefit of real data to minimize the sim-to-real gap; it leverages the scalability of simulation; and it can generalize a pretrained VLA to a target domain with fully automated versus while deep learning on large and diverse datasets has shown promise as a path towards such generalization in computer vision and natural language, collecting high quality datasets of robotic interaction at scale remains an open challenge \citep{Fang2025Rebot,Chen2021Learning,Yu2023Scaling}.

A persistent caveat is that generalist VLAs fail to acquire the task entirely, even under standard in-distribution conditions \citep{Mazza2026Improving}.

For example, domain weights learned by Re-Mix outperform uniform weights by 38\textbackslash{}\% on average and outperform human-selected weights by 32\textbackslash{}\% on datasets used to train existing generalist robot policies, specifically the RT-X models \citep{Hejna2024Optimizing}.

A persistent caveat is that embodiment mixture: naively pooling heterogeneous robot datasets often induces negative transfer rather than gains, underscoring the fragility of indiscriminate data scaling \citep{Wang2026Rethinking}.

\section{Evaluation, Safety \& Deployment}


Evaluation, Safety \& Deployment is where differences in evaluation and deployment and robustness and failure analysis stop looking like implementation detail and start changing what the evidence can support \citep{Kim2026Molmospaces,Francis2021Core,Li2026Robointer,Xu2025Anatomy}.

Across Benchmarks, metrics, and generalization, Safety, reliability, and real-world deployment, and Failure analysis, robustness, and stress testing, the same problem keeps returning: evaluation and deployment and robustness and failure analysis alter the interpretation of otherwise similar results \citep{Luo2026Being,Makarova2025Conther,Liu2026Rdt2,Chen2024Gravmad}.

\subsection{Benchmarks, metrics, and generalization}


A useful way to read benchmarks, metrics, and generalization is to ask which claims survive a change in assumptions: A tension around benchmark coverage and task realism and metric validity and success criteria, motivating a protocol-aware synthesis rather than per-paper summaries. \citep{Liu2026Adaptation,Liu2026World,Maddukuri2025Real,Moroncelli2024Duality}. That question becomes unavoidable when benchmark success versus deployment realism: strong results in curated settings do not automatically translate into safe, reliable embodied behavior under real-world constraints..

A recurrent split along benchmark coverage and task realism runs between planning-centric systems and memory-augmented systems: the former is usually supported by results showing that metrics: simulated environments, as well as the datasets used for EVLP tasks, whereas the latter is more often justified by evidence that many robot learning methods treat such datasets as multi-task expert data and learn a multi-task, generalist policy by training broadly across them \citep{Francis2021Core,Francis2023Core,Memmel2024Strap}. One concrete result is that empirical data confirm the superiority of the algorithm by an average of 38.46\% over other considered methods, and the most successful baseline by 28.21\%, showing higher success rates and faster convergence in the point-reaching task \citep{Makarova2025Conther}. These comparisons are usually reported against mLLMs, EVLP, VLA, GPU, VLM, VLA-oriented, LingBot-VLA, STRAP, CONTHER, HER \citep{Mei2026Video,Zheng2025Multimodal,Francis2023Core,Francis2021Core}.

Seen through benchmark coverage and task realism, the literature separates planning-centric systems from transformer-based policies less by labels alone than by the evidence each side stresses: metrics: simulated environments, as well as the datasets used for EVLP tasks versus empirical data confirm the superiority of the algorithm by an average of 38.46\% over other considered methods, and the most successful baseline by 28.21\%, showing higher success rates and faster convergence in the point-reaching task \citep{Francis2021Core,Francis2023Core,Makarova2025Conther}. Representative evidence comes from reports that a benchmark suite of 8 tasks in which robots interact with the diverse scenes and richly annotated objects \citep{Kim2026Molmospaces}.

The comparison becomes weaker once their impressive capabilities address many of the long-standing challenges faced by physics-based simulators, driving broad adoption in many problem domains, e.g., robotics \citep{Mei2026Video}. A recurrent split along benchmark coverage and task realism runs between memory-augmented systems and transformer-based policies: the former is usually supported by results showing that many robot learning methods treat such datasets as multi-task expert data and learn a multi-task, generalist policy by training broadly across them, whereas the latter is more often justified by evidence that empirical data confirm the superiority of the algorithm by an average of 38.46\% over other considered methods, and the most successful baseline by 28.21\%, showing higher success rates and faster convergence in the point-reaching task \citep{Memmel2024Strap,Makarova2025Conther}.

Representative evidence comes from reports that the main contribution is a detailed breakdown of the five biggest challenges in: Representation, Execution, Generalization, Safety, and Dataset and Evaluation \citep{Xu2025Anatomy}. That picture has to be read cautiously because video models enable photorealistic, physically consistent deformable-body simulation without making prohibitive simplifying assumptions, which is a major bottleneck in physics-based simulation \citep{Mei2026Video}.

A recurrent split along metric validity and success criteria runs between planning-centric systems and memory-augmented systems: the former is usually supported by results showing that metrics: simulated environments, as well as the datasets used for EVLP tasks, whereas the latter is more often justified by evidence that many robot learning methods treat such datasets as multi-task expert data and learn a multi-task, generalist policy by training broadly across them \citep{Francis2021Core,Francis2023Core,Memmel2024Strap}. That picture has to be read cautiously because they thus overcome the limited expressiveness of language-only abstractions in describing intricate physical interactions \citep{Mei2026Video}.

Seen through metric validity and success criteria, the literature separates planning-centric systems from transformer-based policies less by labels alone than by the evidence each side stresses: metrics: simulated environments, as well as the datasets used for EVLP tasks versus empirical data confirm the superiority of the algorithm by an average of 38.46\% over other considered methods, and the most successful baseline by 28.21\%, showing higher success rates and faster convergence in the point-reaching task \citep{Francis2021Core,Francis2023Core,Makarova2025Conther}. Seen through metric validity and success criteria, the literature separates memory-augmented systems from transformer-based policies less by labels alone than by the evidence each side stresses: many robot learning methods treat such datasets as multi-task expert data and learn a multi-task, generalist policy by training broadly across them versus empirical data confirm the superiority of the algorithm by an average of 38.46\% over other considered methods, and the most successful baseline by 28.21\%, showing higher success rates and faster convergence in the point-reaching task \citep{Memmel2024Strap,Makarova2025Conther}.

For example, empirical evaluations on RLBench show that GravMAD significantly outperforms state-of-the-art methods, with a 28.63\% improvement on novel tasks and a 13.36\% gain on tasks encountered during training \citep{Chen2024Gravmad}. A persistent caveat is that systematic reviews and publicly available benchmarks for these models remain limited \citep{Zheng2025Multimodal}.

Along generalization across embodiments and environments, papers centered on planning-centric systems emphasize that systematic reviews and publicly available benchmarks for these models remain limited, while work on memory-augmented systems more often argues that many robot learning methods treat such datasets as multi-task expert data and learn a multi-task, generalist policy by training broadly across them \citep{Zheng2025Multimodal,Francis2021Core,Memmel2024Strap}. That picture has to be read cautiously because even when combinations of these topics are considered, more focus is placed on describing, e.g., current architectural methods, as opposed to also illustrating high-level challenges and opportunities for the field \citep{Francis2023Core}.

For example, such as language-level specification of constraints and multimodal grounding of safety signals, and the amplified risks and evaluation challenges they introduce \citep{Zhang2025Safe}. For example, systematic reviews and publicly available benchmarks for these models remain limited \citep{Zheng2025Multimodal}.

\subsection{Safety, reliability, and real-world deployment}


Results in safety, reliability, and real-world deployment only stay comparable when they are read against GNN-LLM, OWMM, VLM, OWMM-VLM, SOTA, GPT-4o, HHYHRHY, OWMM-Agent, WMs, LingBot-VA. A tension around benchmark realism and task coverage and sim-to-real and deployment constraints, motivating a protocol-aware synthesis rather than per-paper summaries. \citep{Fang2025From,Fang2025Saga,Field2025Text2Touch,Francis2021Core}.

Along benchmark realism and task coverage, papers centered on agent-style systems emphasize that three key findings mapped to these axes: hierarchical memory systems (Capability axis) are important for long-horizon spatial tasks, while work on control and conditioning schemes more often argues that the model features three carefully crafted designs: a shared latent space, integrating vision and action tokens, driven by a Mixture-of-Transformers (MoT) architecture, a closed-loop rollout mechanism, allowing for ongoing \citep{Felicia2026From,Li2026Causal,Brohan2022Robotics}. For example, three key findings mapped to these axes: hierarchical memory systems (Capability axis) are important for long-horizon spatial tasks \citep{Felicia2026From}. These comparisons are usually reported against GNN-LLM, OWMM, VLM, OWMM-VLM, SOTA, GPT-4o, HHYHRHY, OWMM-Agent, WMs, LingBot-VA \citep{Felicia2026From,Chen2025Owmm,Baraldi2025Safety,Li2026Causal}.

Seen through benchmark realism and task coverage, the literature separates agent-style systems from video-conditioned systems less by labels alone than by the evidence each side stresses: three key findings mapped to these axes: hierarchical memory systems (Capability axis) are important for long-horizon spatial tasks versus the approach has several advantages: it enjoys the benefit of real data to minimize the sim-to-real gap; it leverages the scalability of simulation; and it can generalize a pretrained VLA to a target domain with fully automated \citep{Felicia2026From,Fang2025Rebot}. Representative evidence comes from reports that the approach has several advantages: it enjoys the benefit of real data to minimize the sim-to-real gap; it leverages the scalability of simulation; and it can generalize a pretrained VLA to a target domain with fully automated \citep{Fang2025Rebot}. Along benchmark realism and task coverage, papers centered on control and conditioning schemes emphasize that the model features three carefully crafted designs: a shared latent space, integrating vision and action tokens, driven by a Mixture-of-Transformers (MoT) architecture, a closed-loop rollout mechanism, allowing for ongoing, while work on video-conditioned systems more often argues that the approach has several advantages: it enjoys the benefit of real data to minimize the sim-to-real gap; it leverages the scalability of simulation; and it can generalize a pretrained VLA to a target domain with fully automated \citep{Li2026Causal,Brohan2022Robotics,Fang2025Rebot}.

For example, for real-world evaluation with a Franka robot, ReBot increased the success rates of Octo by 17\% and OpenVLA by 20\% \citep{Fang2025Rebot}. A persistent caveat is that a second challenge is the hallucination from domain shift \citep{Chen2025Owmm}. A recurrent split along sim-to-real and deployment constraints runs between agent-style systems and control and conditioning schemes: the former is usually supported by results showing that three key findings mapped to these axes: hierarchical memory systems (Capability axis) are important for long-horizon spatial tasks, whereas the latter is more often justified by evidence that where it shows significant promise in long-horizon manipulation, data efficiency in post-training, and strong generalizability to novel configurations \citep{Felicia2026From,Li2026Causal}.

The comparison becomes weaker once world Models (WMs) have been introduced to provide embodied agents with the abilities to anticipate future environmental states and fill in knowledge gaps, thereby enhancing agents' ability to plan and execute actions \citep{Baraldi2025Safety}. Seen through sim-to-real and deployment constraints, the literature separates agent-style systems from video-conditioned systems less by labels alone than by the evidence each side stresses: three key findings mapped to these axes: hierarchical memory systems (Capability axis) are important for long-horizon spatial tasks versus the approach has several advantages: it enjoys the benefit of real data to minimize the sim-to-real gap; it leverages the scalability of simulation; and it can generalize a pretrained VLA to a target domain with fully automated \citep{Felicia2026From,Fang2025Rebot}.

One concrete result is that providing insights into their strengths and limitations across different task categories \citep{Zhao2025Humanoid}. A persistent caveat is that yet: despite rapid advances in learning-based planning and control, reliable autonomy in real ocean environments remains fundamentally constrained by tightly coupled physical limits \citep{Xu2026Underwater}.

A recurrent split along sim-to-real and deployment constraints runs between control and conditioning schemes and video-conditioned systems: the former is usually supported by results showing that where it shows significant promise in long-horizon manipulation, data efficiency in post-training, and strong generalizability to novel configurations, whereas the latter is more often justified by evidence that the approach has several advantages: it enjoys the benefit of real data to minimize the sim-to-real gap; it leverages the scalability of simulation; and it can generalize a pretrained VLA to a target domain with fully automated \citep{Li2026Causal,Fang2025Rebot}. Representative evidence comes from reports that using two domains--a robot arm and a humanoid--across diverse tasks, simulation data can enhance real-world task performance by an average of 38\%, even with notable differences between the simulation and real-world data \citep{Maddukuri2025Real}.

The comparison becomes weaker once hydrodynamic uncertainty, partial observability, bandwidth-limited communication, and energy scarcity are not independent challenges; they interact within the closed perception-planning-control loop and often amplify one another over time \citep{Xu2026Underwater}. Seen through safety and reliability failure modes, the literature separates agent-style systems from control and conditioning schemes less by labels alone than by the evidence each side stresses: three key findings mapped to these axes: hierarchical memory systems (Capability axis) are important for long-horizon spatial tasks versus the model features three carefully crafted designs: a shared latent space, integrating vision and action tokens, driven by a Mixture-of-Transformers (MoT) architecture, a closed-loop rollout mechanism, allowing for ongoing \citep{Felicia2026From,Li2026Causal,Brohan2022Robotics}.

That picture has to be read cautiously because modeling these world dynamics, especially for dexterous robotics tasks, poses significant challenges due to limited data coverage and scarce action labels \citep{Gao2026Dreamdojo}. That picture has to be read cautiously because systematic evaluation on multiple challenging out-of-distribution (OOD) benchmarks verifies the significance of the method for simulating open-world, contact-rich tasks, paving the way for general-purpose robot world models \citep{Gao2026Dreamdojo}.

One concrete result is that all demonstrations are gathered using a consistent robotic embodiment, paired with precise subtask-level language annotations to facilitate both training and evaluation \citep{Jiang2025Galaxea}. Representative evidence comes from reports that evaluation in simulation offers a scalable complement to real world evaluations, but the visual and physical domain gap between existing simulation benchmarks and the real world has made them an unreliable signal for policy improvement \citep{Jain2025Polaris}.

\subsection{Failure analysis, robustness, and stress testing}


Results in failure analysis, robustness, and stress testing only stay comparable when they are read against RIG, EVLP, WBC, BFMs, BFM, FMRs, ACT, VLA, VLAs, CRT. For Failure analysis, robustness, and stress testing, failure taxonomy and stress scenarios and robustness criteria under deployment shift is a recurring axis of variation, and results are easiest to interpret when protocols and failure assumptions are \citep{Hundt2022Robots,Jain2025Polaris,Jiang2025Behavior,Jiang2025Galaxea}.

Along failure taxonomy and stress scenarios, papers centered on planning-centric systems emphasize that metrics: simulated environments, as well as the datasets used for EVLP tasks, while work on control and conditioning schemes more often argues that pre-incident phase, and post-incident phase \citep{Zhao2025Synergizing,Francis2021Core,Kojima2025Comprehensive,Yuan2025Survey}. These comparisons are usually reported against RIG, EVLP, WBC, BFMs, BFM, FMRs, ACT, VLA, VLAs, CRT \citep{Zhao2025Synergizing,Francis2023Core,Francis2021Core,Yuan2025Survey}. A persistent caveat is that yet: prior work either incorporates only one of these abilities in an end-to-end agent or integrates multiple specialized models into an agent system, limiting the learning efficiency and generalization of the policy \citep{Zhao2025Synergizing}.

Along failure taxonomy and stress scenarios, papers centered on planning-centric systems emphasize that metrics: simulated environments, as well as the datasets used for EVLP tasks, while work on other mapped systems more often argues that evaluations on three representative VLA architectures, including Pi0, Pi0.5 and OpenVLA OFT, show that these models frequently succeed at tasks despite logically impossible instructions, revealing a strong visual bias in action generation \citep{Zhao2025Synergizing,Francis2021Core,Zhang2026Restoring,Liang2025Bootstrap}. For example, pre-incident phase, and post-incident phase \citep{Kojima2025Comprehensive}. That picture has to be read cautiously because even when combinations of these topics are considered, more focus is placed on describing, e.g., current architectural methods, as opposed to also illustrating high-level challenges and opportunities for the field \citep{Francis2023Core}.

A recurrent split along failure taxonomy and stress scenarios runs between control and conditioning schemes and other mapped systems: the former is usually supported by results showing that pre-incident phase, and post-incident phase, whereas the latter is more often justified by evidence that evaluations on three representative VLA architectures, including Pi0, Pi0.5 and OpenVLA OFT, show that these models frequently succeed at tasks despite logically impossible instructions, revealing a strong visual bias in action generation \citep{Kojima2025Comprehensive,Yuan2025Survey,Zhang2026Restoring,Liang2025Bootstrap}. Representative evidence comes from reports that evaluations on three representative VLA architectures, including Pi0, Pi0.5 and OpenVLA OFT, show that these models frequently succeed at tasks despite logically impossible instructions, revealing a strong visual bias in action generation \citep{Zhang2026Restoring}.

The comparison becomes weaker once while learning-based controllers have shown promise for complex tasks, their reliance on labor-intensive and costly retraining for new scenarios limits real-world applicability \citep{Yuan2025Survey}. Seen through robustness criteria under deployment shift, the literature separates planning-centric systems from control and conditioning schemes less by labels alone than by the evidence each side stresses: metrics: simulated environments, as well as the datasets used for EVLP tasks versus pre-incident phase, and post-incident phase \citep{Zhao2025Synergizing,Francis2021Core,Kojima2025Comprehensive,Yuan2025Survey}.

For example, when combined with Diffusion Policy, AFRO substantially increases manipulation success rates across 16 simulated and 4 real-world tasks, outperforming existing pre-training approaches \citep{Liang2025Bootstrap}. The comparison becomes weaker once current limitations, urgent challenges, and future opportunities, positioning BFMs as a key approach toward scalable and general-purpose humanoid intelligence \citep{Yuan2025Survey}.

Seen through robustness criteria under deployment shift, the literature separates planning-centric systems from other mapped systems less by labels alone than by the evidence each side stresses: metrics: simulated environments, as well as the datasets used for EVLP tasks versus evaluations on three representative VLA architectures, including Pi0, Pi0.5 and OpenVLA OFT, show that these models frequently succeed at tasks despite logically impossible instructions, revealing a strong visual bias in action generation \citep{Zhao2025Synergizing,Francis2021Core,Zhang2026Restoring,Liang2025Bootstrap}. For example, from a detailed search of 3850 articles, systematic extraction and evaluation was used to identify and explore 310 papers in depth \citep{Robinson2023Robotic}.

The comparison becomes weaker once there is much room to study (i) pre-incident risk mitigation strategies, (ii) research that assumes physical interaction with humans, and (iii) essential issues of foundation models themselves \citep{Kojima2025Comprehensive}. Along robustness criteria under deployment shift, papers centered on control and conditioning schemes emphasize that pre-incident phase, and post-incident phase, while work on other mapped systems more often argues that evaluations on three representative VLA architectures, including Pi0, Pi0.5 and OpenVLA OFT, show that these models frequently succeed at tasks despite logically impossible instructions, revealing a strong visual bias in action generation \citep{Kojima2025Comprehensive,Yuan2025Survey,Zhang2026Restoring,Liang2025Bootstrap}.

That picture has to be read cautiously because generalist VLAs fail to acquire the task entirely, even under standard in-distribution conditions \citep{Mazza2026Improving}. A recurrent split along stress-test protocol and recovery behavior runs between planning-centric systems and control and conditioning schemes: the former is usually supported by results showing that metrics: simulated environments, as well as the datasets used for EVLP tasks, whereas the latter is more often justified by evidence that pre-incident phase, and post-incident phase \citep{Zhao2025Synergizing,Francis2021Core,Kojima2025Comprehensive,Yuan2025Survey}.

That picture has to be read cautiously because despite their success in controlled environments, reliable real-world deployment is severely hindered by their fragility to visual disturbances \citep{Orjuela2026Improving}. One concrete result is that in this thesis, I will discuss selected works on two domains: tactile sensing and rehabilitation robots, which are exemplars of data sparsity and scarcity, respectively \citep{Xu2025Robot}.

\section{Discussion}


Across current approaches, the strongest pattern is not a single winning method but a recurring dependence on how design choices and evaluation protocols interact \citep{Fan2025Interleave,Shao2025Large,Gundawar2025Bench,Kim2026Molmospaces}.

That interaction explains why apparently similar methods can report different outcomes under different assumptions, resource budgets, or evaluation setups \citep{Din2025Vision,Zheng2024Survey,Zhu2025Objectvla,Doshi2024Scaling}.

It also suggests that future work should treat evaluation design, protocol transparency, and practical constraints as first-class comparative variables rather than as appendix details \citep{Bu2025Univla,Bai2025Towards,Zhou2025Autoeval,Dai2026Conla}.

\section{Conclusion}


Across this literature, performance is best read as evidence-bounded behavior shaped by coordinated choices about methodology, evaluation, and deployment setting \citep{Jiang2025Behavior,Eze2024Learning,Kachaev2025Mind,Xu2025Aerial}.

A survey that compares approaches without making protocol assumptions explicit will overstate some gains and understate some risks \citep{Wu2024Discrete,Lin2025Echovla,Shi2025Diversity,Mandi2022Cacti}.

By structuring the literature around shared lenses and comparable evidence, this comparative framing aims to make future surveys more reproducible, more stable, and easier to extend \citep{Wu2025Moto,Awais2023Foundational,Bai2025Embodied,Black2024Vision}.

\appendix

\section{Tables}



\begin{table}[t]
\centering
\scriptsize
\setlength{\tabcolsep}{2.5pt}
\renewcommand{\arraystretch}{1.0}
\caption{Cross-section comparison map for the survey areas.}
\label{tab:a1}
\begin{tabularx}{\linewidth}{YYYl}
\toprule
Survey area & Comparison focus & Representative evidence & Key refs \\
\midrule
3.1 Manipulation, navigation, and task scope & task setting and embodiment assumptions, observation  and  action interface & Planning  and  reasoning loops vs Video  and  temporal generation on task setting and embodiment assumptions\newline Extensive experiments on the AerialVLN and OpenFly benchmark validate the effectiveness of our & \citep{Xu2025Aerial,Kachaev2025Mind,Dai2026Conla,Fan2025Interleave} \\
3.2 Observation, action, and embodiment interfaces & policy backbone and action representation, language  and  vision grounding into robot actions & Transformer-based generators vs Recent representative works on policy backbone and action representation\newline Evaluated across 6 simulations as well as 3 real-world robots, our 0.9B instantiation-X-VLA-0.9B & \citep{Zheng2025Soft,Hou2025Dita,Lee2026Learning,Zhi2026Motif} \\
\bottomrule
\end{tabularx}
\end{table}

\clearpage

\begin{table}[t]
\centering
\scriptsize
\setlength{\tabcolsep}{2.5pt}
\renewcommand{\arraystretch}{1.0}
\caption{Cross-section comparison map for the survey areas. (continued)}
\begin{tabularx}{\linewidth}{YYYl}
\toprule
Survey area & Comparison focus & Representative evidence & Key refs \\
\midrule
3.3 Task diversity, environment realism, and transfer constraints & task setting and embodiment assumptions, observation  and  action interface & Evaluation  and  benchmark-focused vs Control  and  conditioning on task setting and embodiment assumptions\newline We propose 1) utilizing high visual-fidelity simulation for improved sim-to-real transfer, 2) & \citep{Yang2025Robot,Tang2025Geomanip,Brohan2023Vision,Dong2026Capturing} \\
4.1 Vision-language-action and policy backbones & policy backbone and action representation, language  and  vision grounding into robot actions & Planning  and  reasoning loops vs Transformer-based generators on policy backbone and action representation\newline Extensive evaluations across natural language processing benchmarks, robotic simulation & \citep{Huang2025Motvla,Zheng2025Soft,Hou2025Dita,Shao2025Large} \\
\bottomrule
\end{tabularx}
\end{table}

\clearpage

\begin{table}[t]
\centering
\scriptsize
\setlength{\tabcolsep}{2.5pt}
\renewcommand{\arraystretch}{1.0}
\caption{Cross-section comparison map for the survey areas. (continued)}
\begin{tabularx}{\linewidth}{YYYl}
\toprule
Survey area & Comparison focus & Representative evidence & Key refs \\
\midrule
4.2 World models, planning, and reasoning & planning  and  world-model role, long-horizon state handling & Planning  and  reasoning loops vs Video  and  temporal generation on planning  and  world-model role\newline Real-world evaluations further show that Lumo-1 surpasses strong baselines across a wide range of & \citep{Tang2025Mind,Zheng2025Multimodal,Zhu2025Unified,Liu2026World} \\
4.3 Transfer across embodiments and viewpoints & transfer target (embodiment  and  viewpoint  and  environment), adaptation mechanism and invariance & Control  and  conditioning vs Agent frameworks  and  architectures on transfer target (embodiment  and  viewpoint  and  environment)\newline Our extensive evaluation (6k evaluation trials) shows that our approach leads to performant robotic & \citep{Brohan2023Vision,Chen2026Accelerating,Cui2026Contact} \\
\bottomrule
\end{tabularx}
\end{table}

\clearpage

\begin{table}[t]
\centering
\scriptsize
\setlength{\tabcolsep}{2.5pt}
\renewcommand{\arraystretch}{1.0}
\caption{Cross-section comparison map for the survey areas. (continued)}
\begin{tabularx}{\linewidth}{YYYl}
\toprule
Survey area & Comparison focus & Representative evidence & Key refs \\
\midrule
5.1 Pretraining data and supervision & data source and supervision, pretraining versus post-training & Video  and  temporal generation vs Control  and  conditioning on data source and supervision\newline Through simulated and real-world experiments, we show that: (1) UWM enables effective pretraining & \citep{Zhu2025Unified,Li2025Mask2Iv,Jlg2025Robot,Cui2026Contact} \\
5.2 Post-training, feedback, and continual improvement & data source and supervision, pretraining versus post-training & Evaluation  and  benchmark-focused vs Control  and  conditioning on data source and supervision\newline While recent work has advanced such general robot policies, their training and evaluation are often & \citep{Dai2025Manitaskgen,Jain2025Polaris,Li2026Causal,Jlg2025Robot} \\
\bottomrule
\end{tabularx}
\end{table}

\clearpage

\begin{table}[t]
\centering
\scriptsize
\setlength{\tabcolsep}{2.5pt}
\renewcommand{\arraystretch}{1.0}
\caption{Cross-section comparison map for the survey areas. (continued)}
\begin{tabularx}{\linewidth}{YYYl}
\toprule
Survey area & Comparison focus & Representative evidence & Key refs \\
\midrule
5.3 Data quality, scaling, and distribution shift & data quality and coverage bias, scaling efficiency versus data fidelity & Video  and  temporal generation vs Recent representative works on data quality and coverage bias\newline Our approach has several advantages: 1) it enjoys the benefit of real data to minimize the & \citep{Fang2025Rebot,Chen2021Learning,Wang2026Rethinking,Luo2026Being} \\
6.1 Benchmarks, metrics, and generalization & benchmark coverage and task realism, metric validity and success criteria & Planning  and  reasoning loops vs Memory  and  retrieval augmentation on benchmark coverage and task realism\newline We propose a taxonomy to unify these tasks and provide an in-depth analysis and comparison of the & \citep{Francis2021Core,Francis2023Core,Memmel2024Strap,Makarova2025Conther} \\
\bottomrule
\end{tabularx}
\end{table}

\clearpage

\begin{table}[t]
\centering
\scriptsize
\setlength{\tabcolsep}{2.5pt}
\renewcommand{\arraystretch}{1.0}
\caption{Cross-section comparison map for the survey areas. (continued)}
\begin{tabularx}{\linewidth}{YYYl}
\toprule
Survey area & Comparison focus & Representative evidence & Key refs \\
\midrule
6.2 Safety, reliability, and real-world deployment & benchmark realism and task coverage, sim-to-real and deployment constraints & Agent frameworks  and  architectures vs Control  and  conditioning on benchmark realism and task coverage\newline Our review is complemented by an empirical analysis, wherein we collect and examine predictions & \citep{Baraldi2025Safety,Li2026Causal,Brohan2022Robotics,Fang2025Rebot} \\
6.3 Failure analysis, robustness, and stress testing & failure taxonomy and stress scenarios, robustness criteria under deployment shift & Planning  and  reasoning loops vs Control  and  conditioning on failure taxonomy and stress scenarios\newline The joint learning of reasoning and next image generation explicitly models the inherent & \citep{Zhao2025Synergizing,Francis2021Core,Orjuela2026Improving,Robinson2023Robotic} \\
\bottomrule
\end{tabularx}
\end{table}



\begin{table}[t]
\centering
\scriptsize
\setlength{\tabcolsep}{2.5pt}
\renewcommand{\arraystretch}{1.0}
\caption{Concrete evaluation anchors and protocol cues by survey area.}
\label{tab:a2}
\begin{tabularx}{\linewidth}{YYl}
\toprule
Survey area & Anchor facts and protocol cues & Key refs \\
\midrule
3.1 Manipulation, navigation, and task scope & Comprehensive evaluation in simulation and the real world shows that Interleave-VLA offers two major benefits: (1)\newline Building on this foundation, we present an in-depth examination of large VLM-based VLA models: (1) integration with\newline We also design MolmoSpaces-Bench, a benchmark suite of 8 tasks in which robots interact with our diverse scenes and & \citep{Fan2025Interleave,Shao2025Large,Kim2026Molmospaces,Li2026Momastage} \\
3.2 Observation, action, and embodiment interfaces & Evaluated across 6 simulations as well as 3 real-world robots, our 0.9B instantiation-X-VLA-0.9B simultaneously\newline Evaluations across both simulation and real-world environments validate the superiority of MOTIF, which significantly\newline InternVLA-M1 employs a two-stage pipeline: (i) spatial grounding pre-training on over 2.3M spatial reasoning data to & \citep{Zheng2025Soft,Zhi2026Motif,Chen2025Internvla,Hou2025Dita} \\
\bottomrule
\end{tabularx}
\end{table}

\clearpage

\begin{table}[t]
\centering
\scriptsize
\setlength{\tabcolsep}{2.5pt}
\renewcommand{\arraystretch}{1.0}
\caption{Concrete evaluation anchors and protocol cues by survey area. (continued)}
\begin{tabularx}{\linewidth}{YYl}
\toprule
Survey area & Anchor facts and protocol cues & Key refs \\
\midrule
3.3 Task diversity, environment realism, and transfer constraints & We propose 1) utilizing high visual-fidelity simulation for improved sim-to-real transfer, 2) evaluating policies by\newline Our extensive evaluation (6k evaluation trials) shows that our approach leads to performant robotic policies and\newline Evaluations across both simulation and real-world environments validate the superiority of MOTIF, which significantly & \citep{Yang2025Robot,Brohan2023Vision,Zhi2026Motif,Nasiriany2026Robocasa365} \\
4.1 Vision-language-action and policy backbones & Evaluated across 6 simulations as well as 3 real-world robots, our 0.9B instantiation-X-VLA-0.9B simultaneously\newline Building on this foundation, we present an in-depth examination of large VLM-based VLA models: (1) integration with\newline Specifically, we introduce a unified taxonomy to systematically organize the disparate efforts in this domain, & \citep{Zheng2025Soft,Shao2025Large,Yu2025Survey,Huang2025Motvla} \\
\bottomrule
\end{tabularx}
\end{table}

\clearpage

\begin{table}[t]
\centering
\scriptsize
\setlength{\tabcolsep}{2.5pt}
\renewcommand{\arraystretch}{1.0}
\caption{Concrete evaluation anchors and protocol cues by survey area. (continued)}
\begin{tabularx}{\linewidth}{YYl}
\toprule
Survey area & Anchor facts and protocol cues & Key refs \\
\midrule
4.2 World models, planning, and reasoning & Real-world evaluations further show that Lumo-1 surpasses strong baselines across a wide range of challenging robotic\newline Through simulated and real-world experiments, we show that: (1) UWM enables effective pretraining on large-scale\newline Evaluation mentions include: VLA, VLMs, VLM, MLLMs, FPS, OOD, DreamDojo, VLA-Loop, SANS, UWM. & \citep{Tang2025Mind,Zhu2025Unified,Xu2026Underwater,Bai2025Embodied} \\
4.3 Transfer across embodiments and viewpoints & Our extensive evaluation (6k evaluation trials) shows that our approach leads to performant robotic policies and\newline Reliance on human supervisors imposes a 1:1 supervision ratio that limits scalability, suffers from operator fatigue\newline We show that by conditioning on contact and iterating via simulation, CAP generalizes to novel environments and & \citep{Brohan2023Vision,Chen2026Accelerating,Cui2026Contact,Ter2025Taxonomy} \\
\bottomrule
\end{tabularx}
\end{table}

\clearpage

\begin{table}[t]
\centering
\scriptsize
\setlength{\tabcolsep}{2.5pt}
\renewcommand{\arraystretch}{1.0}
\caption{Concrete evaluation anchors and protocol cues by survey area. (continued)}
\begin{tabularx}{\linewidth}{YYl}
\toprule
Survey area & Anchor facts and protocol cues & Key refs \\
\midrule
5.1 Pretraining data and supervision & Through simulated and real-world experiments, we show that: (1) UWM enables effective pretraining on large-scale\newline We show that by conditioning on contact and iterating via simulation, CAP generalizes to novel environments and\newline Through a comprehensive evaluation on a total of 22 tasks across two robotic manipulation benchmarks (RLBench and & \citep{Zhu2025Unified,Cui2026Contact,Mu2024Adademo,Gao2026Dreamdojo} \\
5.2 Post-training, feedback, and continual improvement & It comprises RoboInter-Tool, a lightweight GUI that enables semi-automatic annotation of diverse representations, and\newline We open-source the code for our semantic autonomous improvement pipeline, as well as our autonomous dataset of 30.5K\newline We evaluate 89 policies over 58,000 simulation rollouts and 2,835 real-world rollouts. & \citep{Li2026Robointer,Zhou2024Autonomous,Lin2026Systematic,Nasiriany2026Robocasa365} \\
\bottomrule
\end{tabularx}
\end{table}

\clearpage

\begin{table}[t]
\centering
\scriptsize
\setlength{\tabcolsep}{2.5pt}
\renewcommand{\arraystretch}{1.0}
\caption{Concrete evaluation anchors and protocol cues by survey area. (continued)}
\begin{tabularx}{\linewidth}{YYl}
\toprule
Survey area & Anchor facts and protocol cues & Key refs \\
\midrule
5.3 Data quality, scaling, and distribution shift & Our approach has several advantages: 1) it enjoys the benefit of real data to minimize the sim-to-real gap; 2) it\newline (2) Embodiment mixture: we find that naively pooling heterogeneous robot datasets often induces negative transfer\newline We empirically demonstrate that Being-H0.5 achieves state-of-the-art results on simulated benchmarks, such as LIBERO & \citep{Fang2025Rebot,Wang2026Rethinking,Luo2026Being,Chen2021Learning} \\
6.1 Benchmarks, metrics, and generalization & Empirical data confirm the superiority of the algorithm by an average of 38.46\% over other considered methods, and the\newline We also design MolmoSpaces-Bench, a benchmark suite of 8 tasks in which robots interact with our diverse scenes and\newline Our main contribution is a detailed breakdown of the five biggest challenges in: (1) Representation, (2) Execution, (3) & \citep{Makarova2025Conther,Kim2026Molmospaces,Xu2025Anatomy,Mei2026Video} \\
\bottomrule
\end{tabularx}
\end{table}

\clearpage

\begin{table}[t]
\centering
\scriptsize
\setlength{\tabcolsep}{2.5pt}
\renewcommand{\arraystretch}{1.0}
\caption{Concrete evaluation anchors and protocol cues by survey area. (continued)}
\begin{tabularx}{\linewidth}{YYl}
\toprule
Survey area & Anchor facts and protocol cues & Key refs \\
\midrule
6.2 Safety, reliability, and real-world deployment & For real-world evaluation with a Franka robot, ReBot increased the success rates of Octo by 17\% and OpenVLA by 20\%.\newline Then, we analyze four main research targets: 1) embodied perception, 2) embodied interaction, 3) embodied agent, and 4)\newline In addition, we conduct an analysis of representative policy learning methods on our dataset, providing insights into & \citep{Fang2025Rebot,Liu2024Aligning,Zhao2025Humanoid,Felicia2026From} \\
6.3 Failure analysis, robustness, and stress testing & The joint learning of reasoning and next image generation explicitly models the inherent correlation between reasoning,\newline In this work, we quantify this vulnerability, demonstrating that state-of-the-art VLAs such as \$π\_\{0.5\}\$ and SmolVLA,\newline From a detailed search of 3850 articles, systematic extraction and evaluation was used to identify and explore 310 & \citep{Zhao2025Synergizing,Orjuela2026Improving,Robinson2023Robotic,Francis2023Core} \\
\bottomrule
\end{tabularx}
\end{table}



\begin{table}[t]
\centering
\scriptsize
\setlength{\tabcolsep}{2.5pt}
\renewcommand{\arraystretch}{1.0}
\caption{Recurring limitations and risk surfaces highlighted across the evidence base.}
\label{tab:a3}
\begin{tabularx}{\linewidth}{YYl}
\toprule
Survey area & Risk or limitation signal & Key refs \\
\midrule
3.1 Manipulation, navigation, and task scope & Indoor mobile manipulation (MoMA) enables robots to translate natural language instructions into physical\newline Learning-based approaches often fail to maintain logical consistency over extended horizons, while methods & \citep{Li2026Momastage} \\
3.2 Observation, action, and embodiment interfaces & While recent vision-language-action models trained on diverse robot datasets exhibit promising generalization\newline However, these intermediate reasoning are often indirect and inherently limited in their capacity to convey & \citep{Hou2025Dita,Zhong2026Acot} \\
\bottomrule
\end{tabularx}
\end{table}

\clearpage

\begin{table}[t]
\centering
\scriptsize
\setlength{\tabcolsep}{2.5pt}
\renewcommand{\arraystretch}{1.0}
\caption{Recurring limitations and risk surfaces highlighted across the evidence base. (continued)}
\begin{tabularx}{\linewidth}{YYl}
\toprule
Survey area & Risk or limitation signal & Key refs \\
\midrule
3.3 Task diversity, environment realism, and transfer constraints & To fill this gap, we present RoboCasa365, a comprehensive simulation benchmark for household mobile\newline By interpreting these constraints through symbolic language representations and translating them into & \citep{Nasiriany2026Robocasa365,Tang2025Geomanip} \\
4.1 Vision-language-action and policy backbones & Despite promising advances, existing approaches face two challenges: limited language steerability when no\newline Pre-trained vision-language-action (VLA) models offer a promising foundation for generalist robot policies, & \citep{Huang2025Motvla,Neary2025Improving} \\
\bottomrule
\end{tabularx}
\end{table}

\clearpage

\begin{table}[t]
\centering
\scriptsize
\setlength{\tabcolsep}{2.5pt}
\renewcommand{\arraystretch}{1.0}
\caption{Recurring limitations and risk surfaces highlighted across the evidence base. (continued)}
\begin{tabularx}{\linewidth}{YYl}
\toprule
Survey area & Risk or limitation signal & Key refs \\
\midrule
4.2 World models, planning, and reasoning & Yet, despite rapid advances in learning-based planning and control, reliable autonomy in real ocean\newline Hydrodynamic uncertainty, partial observability, bandwidth-limited communication, and energy scarcity are not & \citep{Xu2026Underwater} \\
4.3 Transfer across embodiments and viewpoints & Reinforcement learning (RL) has become a foundational approach for enabling intelligent robotic behavior in\newline Equivariant neural networks overcome these limitations by explicitly integrating symmetry and invariance into & \citep{Ter2025Taxonomy,Seo2025Equivariant} \\
\bottomrule
\end{tabularx}
\end{table}

\clearpage

\begin{table}[t]
\centering
\scriptsize
\setlength{\tabcolsep}{2.5pt}
\renewcommand{\arraystretch}{1.0}
\caption{Recurring limitations and risk surfaces highlighted across the evidence base. (continued)}
\begin{tabularx}{\linewidth}{YYl}
\toprule
Survey area & Risk or limitation signal & Key refs \\
\midrule
5.1 Pretraining data and supervision & However, modeling these world dynamics, especially for dexterous robotics tasks, poses significant challenges\newline Systematic evaluation on multiple challenging out-of-distribution (OOD) benchmarks verifies the significance & \citep{Gao2026Dreamdojo} \\
5.2 Post-training, feedback, and continual improvement & To fill this gap, we present RoboCasa365, a comprehensive simulation benchmark for household mobile\newline While recent work has advanced such general robot policies, their training and evaluation are often limited & \citep{Nasiriany2026Robocasa365,Dai2025Manitaskgen} \\
\bottomrule
\end{tabularx}
\end{table}

\clearpage

\begin{table}[t]
\centering
\scriptsize
\setlength{\tabcolsep}{2.5pt}
\renewcommand{\arraystretch}{1.0}
\caption{Recurring limitations and risk surfaces highlighted across the evidence base. (continued)}
\begin{tabularx}{\linewidth}{YYl}
\toprule
Survey area & Risk or limitation signal & Key refs \\
\midrule
5.3 Data quality, scaling, and distribution shift & However, the high cost of real-world data collection hinders further data scaling, thereby restricting the\newline Specifically, ReBot replays real-world robot trajectories in simulation to diversify manipulated objects & \citep{Fang2025Rebot} \\
6.1 Benchmarks, metrics, and generalization & Further, we highlight important challenges hindering the trustworthy integration of video models in robotics,\newline Their impressive capabilities address many of the long-standing challenges faced by physics-based simulators, & \citep{Mei2026Video} \\
\bottomrule
\end{tabularx}
\end{table}

\clearpage

\begin{table}[t]
\centering
\scriptsize
\setlength{\tabcolsep}{2.5pt}
\renewcommand{\arraystretch}{1.0}
\caption{Recurring limitations and risk surfaces highlighted across the evidence base. (continued)}
\begin{tabularx}{\linewidth}{YYl}
\toprule
Survey area & Risk or limitation signal & Key refs \\
\midrule
6.2 Safety, reliability, and real-world deployment & This paper bridges that gap.\newline We conclude by identifying six grand challenges and outlining directions for future research, including the & \citep{Felicia2026From} \\
6.3 Failure analysis, robustness, and stress testing & Yet, prior work either incorporates only one of these abilities in an end-to-end agent or integrates multiple\newline Moreover, even when combinations of these topics are considered, more focus is placed on describing, e.g., & \citep{Zhao2025Synergizing,Francis2023Core} \\
\bottomrule
\end{tabularx}
\end{table}

\bibliographystyle{plainnat}
\bibliography{../citations/ref}

\end{document}
